% ===========================================================================
% CHƯƠNG 3: TRIỂN KHAI THỰC NGHIỆM VÀ ĐÁNH GIÁ KẾT QUẢ
% ===========================================================================
\chapter{TRIỂN KHAI THỰC NGHIỆM VÀ ĐÁNH GIÁ KẾT QUẢ}
% Tiêu đề chạy rút gọn: tiêu đề chương đầy đủ dài hơn một dòng đầu trang,
% khiến fancyhdr phải nới \headheight và làm lệch khung chữ của cả chương.
\markboth{Chương 3. Triển khai thực nghiệm và đánh giá kết quả}{}
\label{chap:thuc-nghiem}

\noindent\textit{Chương 3 trình bày kịch bản thực nghiệm, kết quả huấn luyện mô
hình học máy và đánh giá hiệu năng mô hình kiểm duyệt tín dụng. Nội dung chương
tập trung mô phỏng vận hành và bàn luận, đối chiếu kết quả thực nghiệm
thu được.}

\section{Thiết lập kịch bản và môi trường thực nghiệm}
\label{sec:moi-truong-thuc-nghiem}

\subsection{Hạ tầng phần cứng, phần mềm và thư viện sử dụng}
\label{subsec:ha-tang}

Toàn bộ thực nghiệm được thực hiện trên một máy tính cá nhân với cấu hình và môi
trường phần mềm trình bày tại Bảng~\ref{tab:environment}.

\begin{table}[htbp]
\centering
\caption{Cấu hình phần cứng và môi trường phần mềm}
\label{tab:environment}
\small
\begin{tabular}{@{}L{4.2cm}L{9.6cm}@{}}
\toprule
\textbf{Thành phần} & \textbf{Thông số} \\
\midrule
Hệ điều hành & Windows 11 (64-bit) \\
Ngôn ngữ lập trình & Python 3.12 \\
Xử lý dữ liệu & pandas, NumPy, PyArrow (định dạng lưu trữ Parquet) \\
Mô hình cơ sở & scikit-learn 1.9 \\
Mô hình boosting & XGBoost, LightGBM, CatBoost \\
Giải thích mô hình & SHAP \\
Dịch vụ REST & FastAPI, Uvicorn \\
Trực quan hóa & Matplotlib, Streamlit \\
\bottomrule
\end{tabular}
\end{table}

Các số đo thời gian trong chương này được thực hiện trên một máy đơn, một tiến
trình, nên phản ánh chi phí tính toán của thuật toán chứ không bao gồm thời gian
truyền qua mạng hay ảnh hưởng của tải đồng thời; chúng có giá trị để \textit{so
sánh tương đối} giữa các mô hình (xem thêm mục~\ref{subsec:han-che-dinh-huong}).

Toàn bộ tham số của thực nghiệm được đặt tập trung trong một file cấu hình duy
nhất và hạt giống ngẫu nhiên được cố định, bảo đảm kết quả tái lập được. Quy
trình thực nghiệm gồm sáu bước, mỗi bước là một chương trình độc lập: chuẩn bị dữ
liệu, phân tích và lựa chọn đặc trưng, tinh chỉnh siêu tham số, huấn luyện mô
hình, mô phỏng quy trình kiểm duyệt, và xuất kết quả.

\subsection{Kịch bản thực nghiệm các mô hình học máy}
\label{subsec:kich-ban}

Thực nghiệm được thiết kế theo nguyên tắc \textbf{so sánh công bằng}: toàn bộ
\NModels{} mô hình được huấn luyện trên cùng một tập dữ liệu, cùng một tập đặc
trưng, cùng cách chia dữ liệu và được đánh giá bằng cùng một bộ chỉ tiêu. Mỗi mô
hình đi qua đúng bốn bước:

\begin{enumerate}[label=(\arabic*)]
  \item Huấn luyện trên tập \code{train\_fit}; các mô hình boosting bổ sung cơ
  chế dừng sớm trên tập \code{train\_es};
  \item Hiệu chuẩn xác suất bằng hồi quy đẳng hướng trên tập \code{valid};
  \item Lựa chọn ngưỡng quyết định theo hàm chi phí trên tập \code{valid};
  \item Đánh giá cuối cùng trên tập \code{test} -- chỉ thực hiện một lần.
\end{enumerate}

Về mặt phương pháp, mọi quyết định liên quan tới mô hình đều dựa trên tập
\code{valid}. Tập \code{test} hoàn toàn không tham gia vào bất kỳ lựa
chọn nào, nhờ đó ước lượng hiệu năng cuối cùng không bị lạc quan.

\subsection{Thiết kế và kiểm chứng tập dữ liệu thực nghiệm}
\label{sec:data-design}

\subsubsection{Cơ sở lựa chọn dữ liệu mô phỏng}

Dữ liệu tín dụng cá nhân thuộc nhóm dữ liệu nhạy cảm, chịu ràng buộc chặt chẽ về
bảo mật. Việc tiếp cận dữ liệu thực của một tổ chức tín dụng cho mục đích nghiên
cứu học thuật đòi hỏi quy trình phê duyệt phức tạp và thường không khả thi trong
khuôn khổ một đề án tốt nghiệp. Mặt khác, các bộ dữ liệu tín dụng công khai
thường có ba hạn chế khiến chúng không phù hợp với đề án này: quy mô nhỏ (vài
nghìn quan sát), không có thông tin thời gian nên không cho phép kiểm định ngoài
thời gian, và đặc biệt là \textit{không có dữ liệu giao dịch} -- nhóm dữ liệu mà
đề án muốn đánh giá đóng góp.

Vì vậy đề án sử dụng dữ liệu mô phỏng được sinh bằng chương trình tự xây dựng.
Cách tiếp cận này có một ưu điểm phương pháp luận quan trọng ngoài việc giải
quyết vấn đề tiếp cận dữ liệu: vì quá trình sinh dữ liệu xác định tường minh xác
suất vỡ nợ thật của từng hồ sơ, ta biết được \textbf{trần Bayes} -- giới hạn lý
thuyết mà không mô hình nào có thể vượt qua. Nhờ đó có thể đánh giá một mô hình
còn cách giới hạn lý thuyết bao xa, điều không thể làm được với dữ liệu thực.

Giới hạn tương ứng của cách tiếp cận này là các chỉ tiêu chất lượng thu được
\textbf{không nên được hiểu là dự báo về hiệu quả trên dữ liệu thực}. Giá trị
của thực nghiệm nằm ở việc so sánh \textit{tương đối} giữa các mô hình trong cùng điều kiện và ở việc
kiểm chứng tính đúng đắn của toàn bộ quy trình từ dữ liệu tới quyết định.

\subsubsection{Mô hình sinh dữ liệu biến tiềm ẩn}

Dữ liệu được sinh theo một mô hình xác suất có cấu trúc, trong đó mỗi khách hàng
$i$ được gán một \textbf{biến tiềm ẩn} $z_i \sim \mathcal{N}(0, 1)$ biểu diễn
mức độ tín nhiệm nội tại -- đại lượng không quan sát được trong thực tế. Từ
$z_i$, chương trình sinh ra ba nhóm dữ liệu quan sát được:

\begin{itemize}
  \item \textbf{Thông tin nhân khẩu và tài chính}: thu nhập, nghề nghiệp, thâm
  niên công việc phụ thuộc $z_i$ với nhiễu;
  \item \textbf{Lịch sử tín dụng tại CIC}: điểm CIC, nhóm nợ, số ngày quá hạn là
  các \textit{phép đo có nhiễu} của $z_i$;
  \item \textbf{Hành vi giao dịch}: tính ổn định dòng tiền, tỷ trọng chi tiêu
  nhóm rủi ro cũng phụ thuộc $z_i$ với nhiễu độc lập.
\end{itemize}

Nhãn vỡ nợ được sinh từ một hàm log-odds phi tuyến của $z_i$ và các đặc trưng đã
tính đúng theo thời điểm:

\begin{equation}
\PD_i^{\text{thật}} = \sigma\big( \eta(z_i, \bm{x}_i) + \beta_0 \big),
\qquad
y_i \sim \text{Bernoulli}\big( \PD_i^{\text{thật}} \big),
\label{eq:datagen}
\end{equation}

trong đó hệ số chặn $\beta_0$ được hiệu chỉnh bằng phương pháp số để đạt tỷ lệ vỡ
nợ mục tiêu.

Cấu trúc này tạo ra một đặc điểm quan trọng: vì CIC và hành vi giao dịch chỉ là
các phép đo \textit{có nhiễu} của $z_i$, mô hình học máy chỉ có thể phục hồi một
phần thông tin về $z_i$. Điều này tạo ra một giới hạn trên tự nhiên cho chất
lượng dự báo -- chính là trần Bayes -- giống như tình huống thực tế, nơi không có
mô hình nào đạt $\AUC = 1$.

Quá trình sinh dữ liệu còn đưa vào ba hiện tượng đặc thù của dữ liệu tín dụng
thực: nhóm khách hàng thiếu lịch sử tín dụng (chiếm \ThinFileRate), giá trị khuyết
ở một số trường khai báo, và hiện tượng trôi dạt theo thời gian với mức độ xác
định trước.

\subsubsection{Quy mô và đặc điểm tập dữ liệu}

Bảng~\ref{tab:dataset} trình bày quy mô các bảng dữ liệu thu được. Tập dữ liệu
gồm \NCustomers{} khách hàng, \NAccounts{} tài khoản, \NTransactions{} giao dịch
và \NTotal{} hồ sơ vay trải trên khoảng thời gian hai năm, với thời gian sinh
\GenSeconds{} giây.

\begin{table}[!htbp]
\centering
\caption{Quy mô các bảng dữ liệu sử dụng trong đề án}
\label{tab:dataset}
\small
\setlength{\tabcolsep}{5.0pt}
\begin{tabular}{@{}L{3.0cm}R{2.4cm}L{7.6cm}@{}}
\toprule
\textbf{Bảng dữ liệu} & \textbf{Số bản ghi} & \textbf{Nội dung} \\
\midrule
Khách hàng & 40.000 & Thông tin nhân khẩu học, nghề nghiệp, thu nhập \\
Tài khoản & 73.023 & Tài khoản thanh toán/tiết kiệm và số dư bình quân \\
Giao dịch & 4.538.127 & Giao dịch tài khoản dùng để tính đặc trưng hành vi \\
Hồ sơ vay & 60.000 & Hồ sơ đề nghị cấp tín dụng \\
Bản tin CIC & 60.000 & Thông tin lịch sử tín dụng từ CIC \\
Nhãn rủi ro & 60.000 & Trạng thái vỡ nợ sau cửa sổ quan sát \\
\bottomrule
\end{tabular}
\end{table}


Hình~\ref{fig:data-overview} mô tả phân bố của một số biến chính trong tập dữ
liệu.

\begin{figure}[!htbp]
\centering
\cffig[0.94\textwidth]{data_overview}
\caption{Phân bố các biến chính trong tập dữ liệu thực nghiệm}
\label{fig:data-overview}
\end{figure}

\subsubsection{Kiểm tra chất lượng dữ liệu}
\label{subsec:kq-chat-luong-du-lieu}

Bộ \NConstraints{} ràng buộc nghiệp vụ được áp dụng cho ba bảng dữ liệu chính,
phát hiện \NViolations{} bản ghi vi phạm. Kết quả này không đáng ngạc nhiên với
dữ liệu mô phỏng, nhưng bản thân nó \textit{không chứng minh được} rằng cơ chế
kiểm tra hoạt động đúng -- một bộ kiểm tra hỏng cũng cho kết quả tương tự.

Để phân biệt hai khả năng, đề án thực hiện phép tiêm lỗi: tạo bản sao dữ liệu, cố
ý đưa vào 1\% lỗi ứng với từng loại ràng buộc, rồi chạy lại cơ chế kiểm tra. Kết
quả tại Bảng~\ref{tab:validity} cho thấy toàn bộ lỗi đã tiêm đều được phát hiện
(tỷ lệ phát hiện \ValidityDetected).

\begin{table}[!htbp]
\centering
\caption{Kiểm chứng cơ chế phát hiện dữ liệu không hợp lệ bằng phép tiêm lỗi}
\label{tab:validity}
\small
\setlength{\tabcolsep}{5.0pt}
\begin{tabular}{@{}L{3.3cm}L{5.2cm}R{1.7cm}R{1.7cm}L{1.9cm}@{}}
\toprule
\textbf{Mã ràng buộc} & \textbf{Nội dung ràng buộc} & \textbf{Số lỗi tiêm vào} & \textbf{Số lỗi phát hiện} & \textbf{Kết quả} \\
\midrule
CUS\_\allowbreak AGE\_\allowbreak RANGE & Tuổi khách hàng phải nằm trong khoảng 18-75 & 400 & 400 & Phát hiện đủ \\
CUS\_\allowbreak INCOME\_\allowbreak POSITIVE & Thu nhập khai báo phải lớn hơn 0 & 400 & 400 & Phát hiện đủ \\
APP\_\allowbreak AMOUNT\_\allowbreak POSITIVE & Số tiền đề nghị vay phải lớn hơn 0 & 600 & 600 & Phát hiện đủ \\
APP\_\allowbreak RATE\_\allowbreak RANGE & Lãi suất phải nằm trong khoảng 0-40\%/năm & 600 & 600 & Phát hiện đủ \\
CIC\_\allowbreak SCORE\_\allowbreak RANGE & Điểm CIC phải nằm trong khoảng 300-850 & 600 & 600 & Phát hiện đủ \\
CIC\_\allowbreak GROUP\_\allowbreak ORDER & Nhóm nợ xấu nhất 12 tháng không thể tốt hơn nhóm nợ hiện tại & 600 & 600 & Phát hiện đủ \\
\bottomrule
\end{tabular}

\vspace{2pt}
{\footnotesize\itshape Phép tiêm lỗi được thực hiện trên một bản sao dữ liệu với tỷ lệ 1\% mỗi loại lỗi; dữ liệu gốc không bị thay đổi.}
\end{table}


Về giá trị khuyết, \NMissingVars{} trên tổng số đặc trưng có dữ liệu trống, với
tỷ lệ cao nhất là \MaxMissingPct{} -- tương ứng nhóm khách hàng thiếu lịch sử tín
dụng (Bảng~\ref{tab:missing}). Về giá trị ngoại lệ, \NOutlierVars{} đặc trưng có
quan sát nằm ngoài ngưỡng $3 \times$IQR (Bảng~\ref{tab:outliers}); các giá trị
này được xử lý bằng phép cắt ngọn thay vì loại bỏ quan sát.

\begin{table}[!htbp]
\centering
\caption{Các đặc trưng có giá trị khuyết và chiến lược xử lý}
\label{tab:missing}
\small
\setlength{\tabcolsep}{5.0pt}
\begin{tabular}{@{}L{4.4cm}L{2.0cm}R{2.4cm}L{5.4cm}@{}}
\toprule
\textbf{Đặc trưng} & \textbf{Nhóm} & \textbf{Tỷ lệ khuyết (\%)} & \textbf{Cách xử lý} \\
\midrule
oldest\_\allowbreak account\_\allowbreak months & Lịch sử tín dụng (CIC) & 11,65 & Quy ước nghiệp vụ + cờ chỉ báo khuyết \\
credit\_\allowbreak utilization & Lịch sử tín dụng (CIC) & 10,93 & Quy ước nghiệp vụ + cờ chỉ báo khuyết \\
cic\_score & Lịch sử tín dụng (CIC) & 9,66 & Quy ước nghiệp vụ + cờ chỉ báo khuyết \\
current\_\allowbreak group & Lịch sử tín dụng (CIC) & 9,66 & Quy ước nghiệp vụ + cờ chỉ báo khuyết \\
total\_\allowbreak outstanding & Lịch sử tín dụng (CIC) & 9,66 & Quy ước nghiệp vụ + cờ chỉ báo khuyết \\
dpd\_\allowbreak max\_\allowbreak 12m & Lịch sử tín dụng (CIC) & 9,66 & Quy ước nghiệp vụ + cờ chỉ báo khuyết \\
worst\_\allowbreak group\_\allowbreak 12m & Lịch sử tín dụng (CIC) & 9,66 & Quy ước nghiệp vụ + cờ chỉ báo khuyết \\
num\_\allowbreak bad\_\allowbreak debt\_\allowbreak institutions & Lịch sử tín dụng (CIC) & 9,66 & Quy ước nghiệp vụ + cờ chỉ báo khuyết \\
monthly\_\allowbreak debt\_\allowbreak payment & Lịch sử tín dụng (CIC) & 9,66 & Quy ước nghiệp vụ + cờ chỉ báo khuyết \\
\bottomrule
\end{tabular}
\end{table}


\begin{table}[!htbp]
\centering
\caption{Các đặc trưng có tỷ lệ giá trị ngoại lệ cao nhất (quy tắc $3\times$IQR)}
\label{tab:outliers}
\small
\setlength{\tabcolsep}{5.0pt}
\begin{tabular}{@{}L{5.2cm}R{3.0cm}R{2.6cm}R{2.6cm}@{}}
\toprule
\textbf{Đặc trưng} & \textbf{Tỷ lệ ngoại lệ (\%)} & \textbf{Phân vị 1\%} & \textbf{Phân vị 99\%} \\
\midrule
utility\_\allowbreak amt\_\allowbreak 7d & 24,75 & 0 & 1,49 \\
ecom\_\allowbreak amt\_\allowbreak 7d & 21,90 & 0 & 3,22 \\
current\_\allowbreak group & 20,84 & 1 & 5 \\
outflow\_\allowbreak ratio\_\allowbreak 30d & 18,67 & 0 & 14447060,00 \\
dpd\_\allowbreak x\_\allowbreak institutions & 17,97 & 0 & 87,20 \\
cash\_\allowbreak wd\_\allowbreak amt\_\allowbreak 7d & 14,86 & 0 & 4,91 \\
cash\_\allowbreak ratio\_\allowbreak 7d & 14,86 & 0 & 1,00 \\
cash\_\allowbreak wd\_\allowbreak cnt\_\allowbreak 7d & 14,85 & 0 & 2 \\
repay\_\allowbreak amt\_\allowbreak 7d & 13,29 & 0 & 8,64 \\
repay\_\allowbreak burden\_\allowbreak 7d & 13,29 & 0 & 6209040,00 \\
\bottomrule
\end{tabular}

\vspace{2pt}
{\footnotesize\itshape Giá trị ngoại lệ được xử lý bằng phép cắt ngọn tại phân vị 1\% và 99\% học từ tập huấn luyện, không loại bỏ quan sát.}
\end{table}


\subsubsection{Kết quả gán nhãn rủi ro}
\label{subsec:kq-gan-nhan}

Bảng~\ref{tab:label} trình bày phân bố ba trạng thái rủi ro sau khi gán nhãn theo
tiêu chuẩn nợ nhóm 3--5 với cửa sổ quan sát 12 tháng.

\begin{table}[!htbp]
\centering
\caption{Phân bố trạng thái rủi ro sau khi gán nhãn}
\label{tab:label}
\small
\setlength{\tabcolsep}{5.0pt}
\begin{tabular}{@{}L{8.0cm}R{2.6cm}R{2.6cm}@{}}
\toprule
\textbf{Trạng thái} & \textbf{Số hồ sơ} & \textbf{Tỷ trọng (\%)} \\
\midrule
Tốt (không phát sinh nợ xấu) & 50.867 & 84,78 \\
Xấu (nợ nhóm 3-5, quá hạn ≥ 90 ngày) & 9.133 & 15,22 \\
\bottomrule
\end{tabular}
\end{table}


Cơ sở định lượng cho lựa chọn cửa sổ 12 tháng được trình bày tại
Bảng~\ref{tab:vintage}: đường cong tỷ lệ nợ xấu tích luỹ cho thấy cửa sổ 12 tháng
ghi nhận được \VintageCoverage{} tổng số trường hợp vỡ nợ, trong khi cửa sổ 6
tháng chỉ ghi nhận được khoảng 56\%. Rút ngắn cửa sổ xuống 6 tháng sẽ khiến gần
một nửa số hồ sơ thực sự xấu bị gán nhãn tốt -- sai lệch đủ lớn để làm hỏng mô
hình.

\begin{table}[!htbp]
\centering
\caption{Tỷ lệ nợ xấu tích luỹ theo số tháng kể từ khi cấp tín dụng}
\label{tab:vintage}
\small
\setlength{\tabcolsep}{5.0pt}
\begin{tabular}{@{}R{3.4cm}R{3.5cm}R{3.5cm}R{3.5cm}@{}}
\toprule
\textbf{Tháng kể từ khi cấp tín dụng} & \textbf{Số hồ sơ phát sinh nợ xấu} & \textbf{Tỷ lệ nợ xấu tích luỹ (\%)} & \textbf{Tỷ lệ ghi nhận được (\%)} \\
\midrule
3 & 1.460 & 2,43 & 16,0 \\
6 & 5.078 & 8,46 & 55,6 \\
9 & 7.632 & 12,72 & 83,6 \\
12 & 9.133 & 15,22 & 100 \\
15 & 9.133 & 15,22 & 100 \\
18 & 9.133 & 15,22 & 100 \\
\bottomrule
\end{tabular}

\vspace{2pt}
{\footnotesize\itshape Cửa sổ quan sát 12 tháng được chọn tại điểm đường cong tích luỹ bắt đầu đi ngang.}
\end{table}


\subsection{Phương pháp phân chia tập dữ liệu thử nghiệm}
\label{subsec:kq-chia-tap}

Kết quả phân chia dữ liệu theo thời gian được trình bày tại
Bảng~\ref{tab:split}.

\begin{table}[!htbp]
\centering
\caption{Phân chia dữ liệu theo thời gian}
\label{tab:split}
\small
\setlength{\tabcolsep}{5.0pt}
\begin{tabular}{@{}L{3.6cm}R{2.0cm}R{1.6cm}L{5.6cm}@{}}
\toprule
\textbf{Tập} & \textbf{Số hồ sơ} & \textbf{Tỷ lệ} & \textbf{Vai trò} \\
\midrule
Huấn luyện (train\_\allowbreak fit) & 35.700 & 59,5\% & Ước lượng tham số mô hình \\
Dừng sớm (train\_es) & 6.300 & 10,5\% & Xác định số vòng lặp tối ưu \\
Kiểm định (valid) & 9.000 & 15,0\% & Hiệu chuẩn xác suất, chọn ngưỡng \\
Kiểm thử (test) & 9.000 & 15,0\% & Đánh giá cuối cùng, dùng một lần \\
\bottomrule
\end{tabular}

\vspace{2pt}
{\footnotesize\itshape Giai đoạn huấn luyện: 2024-07-01 đến 2025-11-24; giai đoạn kiểm thử: 2026-03-12 đến 2026-06-30.}
\end{table}


Tỷ lệ nợ xấu tăng dần qua các giai đoạn, từ \DRTrain{}
ở tập huấn luyện lên \DRTest{} ở tập kiểm thử. Đây chính là biểu hiện của hiện
tượng trôi dạt đã được đưa vào quá trình sinh dữ liệu. Sự chênh lệch này khiến
bài toán trở nên khó hơn so với chia ngẫu nhiên và làm cho kết quả đánh giá sát
với điều kiện vận hành thực tế hơn.

\section{Kết quả huấn luyện và tính toán mô hình}
\label{sec:ket-qua-huan-luyen}

\subsection{Kết quả xây dựng và lựa chọn đặc trưng}
\label{subsec:kq-dac-trung}

Tập đặc trưng ban đầu gồm \NFeaturesInitial{} biến, phân bố theo năm nhóm nguồn
dữ liệu như trình bày tại Bảng~\ref{tab:feature-groups}.

\begin{table}[!htbp]
\centering
\caption{Cơ cấu tập đặc trưng theo nhóm và nguồn dữ liệu}
\label{tab:feature-groups}
\small
\setlength{\tabcolsep}{5.0pt}
\begin{tabular}{@{}L{4.2cm}L{4.2cm}R{3.0cm}@{}}
\toprule
\textbf{Nhóm} & \textbf{Nguồn dữ liệu} & \textbf{Số đặc trưng} \\
\midrule
Hành vi giao dịch & GIAO DỊCH & 99 \\
Lịch sử tín dụng (CIC) & BẢN TIN CIC & 16 \\
Thông tin cá nhân & HỒ SƠ KH & 15 \\
Biến tổ hợp & TÍNH TOÁN & 13 \\
Tham số khoản vay & ĐƠN VAY & 11 \\
\bottomrule
\end{tabular}
\end{table}


\subsubsection{Phân tích Information Value}

Bảng~\ref{tab:iv-top} liệt kê mười lăm đặc trưng có năng lực dự báo cao nhất.

\begin{table}[!htbp]
\centering
\caption{Mười lăm đặc trưng có năng lực dự báo cao nhất theo Information Value}
\label{tab:iv-top}
\small
\setlength{\tabcolsep}{5.0pt}
\begin{tabular}{@{}L{5.0cm}L{2.4cm}R{1.8cm}L{5.0cm}@{}}
\toprule
\textbf{Đặc trưng} & \textbf{Nhóm} & \textbf{IV} & \textbf{Mức độ} \\
\midrule
cic\_score & Lịch sử tín dụng (CIC) & 1,244 & Rất mạnh (cần kiểm tra rò rỉ nhãn) \\
worst\_\allowbreak group\_\allowbreak 12m & Lịch sử tín dụng (CIC) & 1,179 & Rất mạnh (cần kiểm tra rò rỉ nhãn) \\
dpd\_\allowbreak max\_\allowbreak 12m & Lịch sử tín dụng (CIC) & 1,106 & Rất mạnh (cần kiểm tra rò rỉ nhãn) \\
current\_\allowbreak group & Lịch sử tín dụng (CIC) & 0,959 & Rất mạnh (cần kiểm tra rò rỉ nhãn) \\
dpd\_\allowbreak x\_\allowbreak institutions & Biến tổ hợp & 0,948 & Rất mạnh (cần kiểm tra rò rỉ nhãn) \\
num\_\allowbreak bad\_\allowbreak debt\_\allowbreak institutions & Lịch sử tín dụng (CIC) & 0,821 & Rất mạnh (cần kiểm tra rò rỉ nhãn) \\
disposable\_\allowbreak income & Biến tổ hợp & 0,739 & Rất mạnh (cần kiểm tra rò rỉ nhãn) \\
disposable\_\allowbreak ratio & Biến tổ hợp & 0,660 & Rất mạnh (cần kiểm tra rò rỉ nhãn) \\
total\_\allowbreak avg\_\allowbreak balance & Thông tin cá nhân & 0,564 & Rất mạnh (cần kiểm tra rò rỉ nhãn) \\
inflow\_90d & Hành vi giao dịch & 0,546 & Rất mạnh (cần kiểm tra rò rỉ nhãn) \\
inflow\_\allowbreak 180d & Hành vi giao dịch & 0,521 & Rất mạnh (cần kiểm tra rò rỉ nhãn) \\
monthly\_\allowbreak income & Thông tin cá nhân & 0,503 & Rất mạnh (cần kiểm tra rò rỉ nhãn) \\
txn\_\allowbreak amt\_\allowbreak std\_\allowbreak 90d & Hành vi giao dịch & 0,448 & Mạnh \\
txn\_\allowbreak amt\_\allowbreak std\_\allowbreak 180d & Hành vi giao dịch & 0,443 & Mạnh \\
net\_\allowbreak flow\_\allowbreak 180d & Hành vi giao dịch & 0,432 & Mạnh \\
\bottomrule
\end{tabular}
\end{table}


Kết quả cho thấy nhóm đặc trưng CIC chiếm ưu thế ở phần đầu bảng xếp hạng, với
giá trị $\IV$ cao nhất đạt \MaxIV{}. Điều này phù hợp với trực giác nghiệp vụ:
lịch sử trả nợ trong quá khứ là chỉ báo mạnh nhất về khả năng trả nợ trong tương
lai. Giá trị $\IV$ vượt ngưỡng 0,5 thông thường là dấu hiệu cần kiểm tra rò rỉ
nhãn; trong trường hợp này, việc kiểm tra ràng buộc thời gian
(CIC\_QUERY\_BEFORE\_APPLY) xác nhận mọi bản tin CIC đều được truy vấn trước thời
điểm nộp hồ sơ, nên $\IV$ cao phản ánh bản chất của dữ liệu chứ không phải lỗi
phương pháp.

Bảng~\ref{tab:iv-group} tổng hợp $\IV$ theo nhóm nguồn dữ liệu.

\begin{table}[!htbp]
\centering
\caption{Information Value tổng hợp theo nhóm nguồn dữ liệu}
\label{tab:iv-group}
\small
\setlength{\tabcolsep}{5.0pt}
\begin{tabular}{@{}L{3.426cm}R{2.474cm}R{2.474cm}R{2.855cm}R{2.665cm}@{}}
\toprule
\textbf{Nhóm} & \textbf{Số đặc trưng} & \textbf{Tổng IV} & \textbf{IV trung bình} & \textbf{IV lớn nhất} \\
\midrule
Hành vi giao dịch & 99 & 13,945 & 0,141 & 0,546 \\
Lịch sử tín dụng (CIC) & 16 & 6,780 & 0,424 & 1,244 \\
Biến tổ hợp & 13 & 4,253 & 0,327 & 0,948 \\
Thông tin cá nhân & 15 & 1,604 & 0,107 & 0,564 \\
Tham số khoản vay & 11 & 0,589 & 0,054 & 0,190 \\
\bottomrule
\end{tabular}
\end{table}


Xét theo $\IV$ trung bình, nhóm CIC dẫn đầu rõ rệt: mỗi đặc trưng CIC mang nhiều
thông tin hơn hẳn. Nhưng xét theo \textit{tổng} $\IV$, nhóm hành vi giao dịch lại
chiếm vị trí cao nhất nhờ số lượng đặc trưng lớn. Hai cách nhìn này không mâu
thuẫn mà bổ sung cho nhau: dữ liệu CIC là nguồn thông tin cô đọng nhất, còn dữ
liệu giao dịch là nguồn thông tin dồi dào nhất. Với nhóm khách hàng thiếu lịch sử
tín dụng, nơi dữ liệu CIC gần như trống, nhóm thứ hai trở thành nguồn thông tin
chủ yếu.

\begin{figure}[!htbp]
\centering
\cffig[0.95\textwidth]{iv_ranking}
\caption{Xếp hạng đặc trưng theo Information Value}
\label{fig:iv-ranking}
\end{figure}

\subsubsection{Kết quả lựa chọn tập đặc trưng}

Quy trình lọc ba bước rút gọn tập đặc trưng từ \NFeaturesInitial{} xuống
\NFeaturesSelected{} biến (Bảng~\ref{tab:selection}).

\begin{table}[!htbp]
\centering
\caption{Kết quả lựa chọn tập đặc trưng qua ba bộ lọc}
\label{tab:selection}
\small
\setlength{\tabcolsep}{5.0pt}
\begin{tabular}{@{}L{7.4cm}R{3.2cm}R{3.2cm}@{}}
\toprule
\textbf{Bước} & \textbf{Số đặc trưng} & \textbf{Đã loại} \\
\midrule
Tập đặc trưng ban đầu & 154 & 0 \\
Sau lọc theo IV (≥ 0.02) & 126 & 28 \\
Sau lọc theo tương quan (|r| ≤ 0.95) & 116 & 10 \\
Sau lọc theo độ ổn định (PSI ≤ 0.25) & 116 & 0 \\
\bottomrule
\end{tabular}
\end{table}


\begin{table}[!htbp]
\centering
\caption{Cơ cấu tập đặc trưng trước và sau lựa chọn}
\label{tab:selection-group}
\small
\setlength{\tabcolsep}{5.0pt}
\begin{tabular}{@{}L{4.0cm}R{2.8cm}R{3.0cm}R{3.2cm}@{}}
\toprule
\textbf{Nhóm} & \textbf{Ban đầu} & \textbf{Sau lựa chọn} & \textbf{Tỷ lệ giữ lại (\%)} \\
\midrule
Hành vi giao dịch & 99 & 76 & 76,8 \\
Lịch sử tín dụng (CIC) & 16 & 14 & 87,5 \\
Thông tin cá nhân & 15 & 8 & 53,3 \\
Biến tổ hợp & 13 & 12 & 92,3 \\
Tham số khoản vay & 11 & 6 & 54,5 \\
\bottomrule
\end{tabular}
\end{table}


Phân tích cơ cấu sau lựa chọn (Bảng~\ref{tab:selection-group}) cho thấy các bộ
lọc hoạt động đúng như thiết kế: nhóm biến tổ hợp giữ lại tỷ lệ cao nhất -- hợp
lý vì mỗi biến trong nhóm này được thiết kế để biểu diễn một chỉ tiêu nghiệp vụ
riêng biệt nên ít trùng lặp. Ngược lại, nhóm tham số khoản vay bị loại nhiều
nhất, chủ yếu do năng lực dự báo thấp.

Bước lọc theo độ ổn định không loại bỏ đặc trưng nào: chỉ số $\PSI$
lớn nhất trên toàn bộ tập đặc trưng chỉ đạt \MaxPSI{}, thấp hơn nhiều so với
ngưỡng cảnh báo 0,10. Kết quả này cho thấy mức trôi dạt trong dữ liệu thực nghiệm
tập trung ở \textit{quan hệ giữa đặc trưng và nhãn} chứ không ở phân phối biên
của đặc trưng -- một dạng trôi dạt khó phát hiện hơn bằng giám sát $\PSI$ đơn
thuần.

\subsection{Kết quả trên các mô hình cơ sở}
\label{subsec:kq-mo-hinh-co-so}

Nhóm mô hình cơ sở gồm hồi quy logistic, thẻ điểm WOE và cây quyết định -- các mô
hình có cấu trúc đơn giản và khả năng giải thích cao.

Hồi quy logistic đạt $\AUC = \LRAUC{}$ trên tập kiểm thử, thẻ điểm WOE đạt
\WOEAUC{}, còn cây quyết định đơn đạt \DTAUC{}. Khoảng cách giữa hai mô hình
tuyến tính và cây quyết định là đáng kể -- gần 0,04 điểm $\AUC$ -- cho thấy hạn
chế về phương sai cao của cây đơn đã được phân tích ở
mục~\ref{subsec:decision-tree} thể hiện rõ trong thực nghiệm.

Việc hồi quy logistic -- mô hình đơn giản nhất trong nhóm -- lại đạt kết quả gần
sát với các mô hình phức tạp là một quan sát quan trọng, sẽ được bàn kỹ ở
mục~\ref{subsec:so-sanh-thuat-toan}.

\subsection{Kết quả trên các mô hình học máy nâng cao}
\label{subsec:kq-mo-hinh-nang-cao}

Nhóm mô hình nâng cao gồm các mô hình học tập hợp (rừng ngẫu nhiên, Extra Trees,
bốn biến thể gradient boosting), máy vectơ hỗ trợ, mạng nơ-ron và mô hình kết hợp
xếp tầng.

Trong nhóm này, mô hình kết hợp xếp tầng đạt $\AUC$ danh nghĩa cao nhất
(\TopAUC{}), ngay sau là ba biến thể gradient boosting: CatBoost \CatAUC{},
XGBoost \XGBAUC{} và LightGBM \LGBAUC{}. Máy vectơ hỗ trợ đạt \SVMAUC{}, rừng
ngẫu nhiên đạt \RFAUC{} và Hist Gradient Boosting đạt \GBAUC{}.

Cần nhấn mạnh ngay rằng khoảng cách giữa các mô hình dẫn đầu là \textit{không
đáng kể}: chênh lệch giữa mô hình đứng đầu và mô hình thứ hai chỉ \GapTopSecond{}
điểm $\AUC$. Mục~\ref{subsec:kq-chon-mo-hinh} phân tích vì sao mức chênh lệch
này không đủ cơ sở để kết luận một thuật toán vượt trội, và vì sao mô hình được
chọn triển khai lại không phải mô hình có $\AUC$ cao nhất.

Kết quả của SVM cần được bình luận riêng. Mặc dù sử dụng phép xấp xỉ Nystroem
thay vì giải bài toán đối ngẫu đầy đủ, mô hình vẫn đạt mức chất lượng tương đương
rừng ngẫu nhiên và chỉ kém mô hình tốt nhất \GapBestSVM{} điểm $\AUC$, với thời
gian huấn luyện khoảng \SVMTrainSec{} giây. Kết quả này cho thấy phép xấp xỉ
hạt nhân là giải pháp khả thi để áp dụng SVM ở quy mô dữ liệu mà SVM hạt nhân đầy đủ không huấn
luyện được trong thời gian chấp nhận được.

\subsection{Tối ưu hóa siêu tham số cho mô hình}
\label{subsec:kq-tuning}

Kết quả trình bày ở hai mục trên đã sử dụng siêu tham số được tinh chỉnh ở bước
này; mục này trình bày chi tiết quá trình tinh chỉnh đó.

Ba mô hình có nhiều siêu tham số nhất -- LightGBM, XGBoost và rừng ngẫu nhiên --
được tinh chỉnh bằng tìm kiếm ngẫu nhiên kết hợp kiểm định chéo theo chuỗi thời
gian trên tập huấn luyện. Kết quả trình bày tại Bảng~\ref{tab:tuning}.

\begin{table}[!htbp]
\centering
\caption{Kết quả tinh chỉnh siêu tham số bằng tìm kiếm ngẫu nhiên}
\label{tab:tuning}
\small
\setlength{\tabcolsep}{5.0pt}
\begin{tabular}{@{}L{4.6cm}R{2.6cm}R{2.6cm}R{3.4cm}@{}}
\toprule
\textbf{Mô hình} & \textbf{Số lần thử} & \textbf{AUC (CV)} & \textbf{Thời gian (phút)} \\
\midrule
lightgbm & 14 & 0,8499 & 1,5 \\
xgboost & 14 & 0,8623 & 2,9 \\
random\_\allowbreak forest & 14 & 0,8576 & 9,5 \\
\bottomrule
\end{tabular}

\vspace{2pt}
{\footnotesize\itshape Kiểm định chéo theo chuỗi thời gian (TimeSeriesSplit) ba lần khớp trên tập huấn luyện.}
\end{table}


Hai điểm cần lưu ý về phương pháp. Thứ nhất, việc tinh chỉnh được thực hiện trên
\textit{đúng tập đặc trưng} mà mô hình cuối cùng sử dụng -- tập \NFeaturesSelected{}
biến sau lựa chọn, không phải tập đầy đủ ban đầu. Nếu tinh chỉnh trên tập đầy đủ
rồi áp dụng tham số cho tập đã rút gọn, siêu tham số tìm được không còn tối ưu.
Thứ hai, toàn bộ quá trình tinh chỉnh chỉ sử dụng tập huấn luyện; tập kiểm định
và tập kiểm thử không tham gia, nhờ đó không tạo ra rò rỉ thông tin.

\section{Đánh giá hiệu năng và so sánh kết quả thực nghiệm}
\label{sec:danh-gia-hieu-nang}

\subsection{Đánh giá qua các chỉ số phân loại}
\label{subsec:kq-chi-so-phan-loai}

Bảng~\ref{tab:model-quality} trình bày đầy đủ kết quả của \NModels{} mô hình theo
bảy chỉ tiêu chất lượng dự báo.

\begin{table}[!htbp]
\centering
\caption{So sánh chất lượng dự báo của các mô hình trên tập kiểm thử}
\label{tab:model-quality}
\small
\setlength{\tabcolsep}{3.0pt}
\begin{tabular}{@{}L{4.374cm}R{1.350cm}R{1.350cm}R{1.350cm}R{1.350cm}R{1.350cm}R{1.350cm}R{1.350cm}@{}}
\toprule
\textbf{Mô hình} & \textbf{AUC} & \textbf{Gini} & \textbf{KS} & \textbf{PR-AUC} & \textbf{F1} & \textbf{Brier} & \textbf{ECE} \\
\midrule
\textbf{Kết hợp xếp tầng (Stacking)} & \textbf{0,8620} & \textbf{0,7240} & \textbf{0,5702} & \textbf{0,6141} & \textbf{0,4445} & \textbf{0,0981} & \textbf{0,0063} \\
CatBoost & 0,8614 & 0,7229 & 0,5689 & 0,6257 & 0,4731 & 0,0976 & 0,0075 \\
XGBoost & 0,8612 & 0,7223 & 0,5729 & 0,6213 & 0,4689 & 0,0980 & 0,0101 \\
LightGBM & 0,8590 & 0,7181 & 0,5647 & 0,6107 & 0,4482 & 0,0988 & 0,0053 \\
Hồi quy logistic & 0,8579 & 0,7158 & 0,5629 & 0,5914 & 0,4418 & 0,1007 & 0,0070 \\
Gradient Boosting (histogram) & 0,8575 & 0,7150 & 0,5594 & 0,5999 & 0,4453 & 0,0993 & 0,0107 \\
Thẻ điểm WOE + Hồi quy logistic & 0,8570 & 0,7140 & 0,5621 & 0,6025 & 0,4627 & 0,1000 & 0,0112 \\
Mạng nơ-ron truyền thẳng (MLP) & 0,8560 & 0,7119 & 0,5602 & 0,6012 & 0,4649 & 0,0999 & 0,0083 \\
Máy vectơ hỗ trợ (SVM-RBF) & 0,8540 & 0,7080 & 0,5619 & 0,5953 & 0,4517 & 0,1005 & 0,0070 \\
Rừng ngẫu nhiên & 0,8527 & 0,7054 & 0,5531 & 0,5930 & 0,4654 & 0,1008 & 0,0064 \\
Cây cực ngẫu nhiên (Extra Trees) & 0,8471 & 0,6942 & 0,5448 & 0,5808 & 0,4360 & 0,1030 & 0,0078 \\
Cây quyết định (CART) & 0,8206 & 0,6412 & 0,4970 & 0,5299 & 0,4052 & 0,1092 & 0,0066 \\
\bottomrule
\end{tabular}

\vspace{2pt}
{\footnotesize\itshape Sắp xếp theo ROC-AUC giảm dần. Trần Bayes (tính từ xác suất vỡ nợ thực của quá trình sinh dữ liệu) là 0,9109.}
\end{table}


Mô hình được chọn triển khai -- \BestModel{} -- đạt $\AUC = \BestAUC{}$,
$\text{Gini} = \BestGini{}$, $\KS = \BestKS{}$ và $F_1 = \BestFOne{}$. Về chất
lượng hiệu chuẩn, điểm Brier đạt \BestBrier{} và sai số hiệu chuẩn kỳ vọng chỉ
\BestECE{}; mức sai số thấp này cho phép sử dụng trực tiếp xác suất đầu ra làm
căn cứ phân vùng quyết định.

Hình~\ref{fig:roc-curves} và Hình~\ref{fig:calibration} minh họa khả năng phân
biệt và mức độ hiệu chuẩn của các mô hình.

\begin{figure}[!htbp]
\centering
\cffig[0.95\textwidth]{roc_curves}
\caption{Đường cong ROC của các mô hình trên tập kiểm thử}
\label{fig:roc-curves}
\end{figure}

\begin{figure}[!htbp]
\centering
\cffig[0.95\textwidth]{calibration}
\caption{Đường cong hiệu chuẩn xác suất sau khi hiệu chuẩn đẳng hướng}
\label{fig:calibration}
\end{figure}

\subsubsection{So sánh với trần Bayes}

Vì quá trình sinh dữ liệu xác định tường minh xác suất vỡ nợ thật của từng hồ sơ,
có thể tính được trần Bayes -- giới hạn lý thuyết của mọi mô hình trên tập dữ
liệu này. Trần Bayes trên tập kiểm thử là \BayesAUC{}.

Mô hình tốt nhất đạt \BestAUC{}, tức còn cách trần Bayes \GapBayes{} điểm
$\AUC$. Khoảng cách này có ý nghĩa quan trọng về mặt diễn giải: nó cho thấy mô
hình đã khai thác được phần lớn thông tin có trong dữ liệu, và phần còn thiếu
không phải do thuật toán yếu mà do bản thân dữ liệu quan sát được chỉ là phép đo
có nhiễu của mức tín nhiệm thật. Nói cách khác, việc đầu tư thêm vào thuật toán
phức tạp hơn sẽ cho lợi ích giảm dần; cải thiện đáng kể hơn chỉ đến từ việc bổ
sung nguồn dữ liệu mới mang thông tin bổ sung.

\subsection{Đánh giá khả năng phân tách rủi ro}
\label{subsec:kq-phan-tach}

Khả năng phân tách rủi ro được đánh giá qua ba công cụ bổ trợ nhau.

\paragraph{Thống kê KS.} Mô hình được chọn đạt $\KS = \BestKS{}$, nghĩa là tại
điểm cắt tối ưu, khoảng cách giữa phân phối tích lũy của nhóm tốt và nhóm xấu đạt
\KSPoints{} điểm phần trăm. Theo thực hành ngành, giá trị $\KS$ trên 0,40 được
coi là mô hình có khả năng phân tách tốt.

\begin{figure}[!htbp]
\centering
\cffig[0.9\textwidth]{ks_curve}
\caption{Đường cong Kolmogorov--Smirnov của mô hình được chọn triển khai}
\label{fig:ks-curve}
\end{figure}

\paragraph{Phân tích theo thập phân vị.} Bảng~\ref{tab:gains} trình bày phân bố
hồ sơ xấu theo mười nhóm điểm rủi ro.

\begin{table}[!htbp]
\centering
\caption{Phân tích theo thập phân vị điểm rủi ro}
\label{tab:gains}
\small
\setlength{\tabcolsep}{3.0pt}
\begin{tabular}{@{}R{1.113cm}R{1.370cm}R{1.455cm}R{1.712cm}R{1.712cm}R{1.798cm}R{1.798cm}R{1.113cm}R{1.541cm}@{}}
\toprule
\textbf{Nhóm} & \textbf{Số hồ sơ} & \textbf{Số hồ sơ xấu} & \textbf{PD trung bình} & \textbf{Tỷ lệ nợ xấu (\%)} & \textbf{Tỷ trọng nợ xấu (\%)} & \textbf{Nợ xấu tích luỹ (\%)} & \textbf{Lift} & \textbf{Lift tích luỹ} \\
\midrule
1 & 900 & 647 & 0,7130 & 71,89 & 40,67 & 40,67 & 4,07 & 4,07 \\
2 & 900 & 362 & 0,3976 & 40,22 & 22,75 & 63,42 & 2,28 & 3,17 \\
3 & 900 & 213 & 0,2196 & 23,67 & 13,39 & 76,81 & 1,34 & 2,56 \\
4 & 900 & 128 & 0,1556 & 14,22 & 8,05 & 84,85 & 0,80 & 2,12 \\
5 & 900 & 85 & 0,0975 & 9,44 & 5,34 & 90,19 & 0,53 & 1,80 \\
6 & 900 & 59 & 0,0663 & 6,56 & 3,71 & 93,90 & 0,37 & 1,57 \\
7 & 900 & 43 & 0,0482 & 4,78 & 2,70 & 96,61 & 0,27 & 1,38 \\
8 & 900 & 21 & 0,0258 & 2,33 & 1,32 & 97,93 & 0,13 & 1,22 \\
9 & 900 & 18 & 0,0119 & 2 & 1,13 & 99,06 & 0,11 & 1,10 \\
10 & 900 & 15 & 0,0091 & 1,67 & 0,94 & 100 & 0,09 & 1 \\
\bottomrule
\end{tabular}

\vspace{2pt}
{\footnotesize\itshape Nhóm 1 gồm 10\% hồ sơ có điểm rủi ro cao nhất. Lift là tỷ số giữa mật độ nợ xấu của nhóm và mật độ nợ xấu trung bình.}
\end{table}


Kết quả cho thấy thập phân vị rủi ro cao nhất tập trung \TopDecileBad{} tổng số
hồ sơ xấu -- gấp hơn bốn lần mức phân bố ngẫu nhiên. Ba thập phân vị rủi ro cao
nhất chứa \TopThreeDecileBad{} tổng số hồ sơ xấu. Đây là thước đo trực tiếp giá trị vận hành của mô hình:
nếu tổ chức tín dụng chỉ có nguồn lực để thẩm định kỹ 30\% số hồ sơ, việc ưu tiên
theo điểm rủi ro của mô hình cho phép bao phủ hơn ba phần tư số hồ sơ có vấn đề.

\begin{figure}[!htbp]
\centering
\cffig[0.95\textwidth]{gains_lift}
\caption{Đường cong gains và lift theo thập phân vị điểm rủi ro}
\label{fig:gains}
\end{figure}

\paragraph{Phân bố điểm rủi ro.} Hình~\ref{fig:score-distribution} cho thấy phân
bố điểm rủi ro của hai nhóm hồ sơ và vị trí của các ngưỡng quyết định.

\begin{figure}[!htbp]
\centering
\cffig[0.95\textwidth]{score_distribution}
\caption{Phân bố xác suất vỡ nợ dự báo theo nhóm hồ sơ thực tế}
\label{fig:score-distribution}
\end{figure}

\subsection{So sánh hiệu năng giữa các thuật toán}
\label{subsec:so-sanh-thuat-toan}

\subsubsection{Chênh lệch giữa các thuật toán là nhỏ}

Quan sát nổi bật nhất từ Bảng~\ref{tab:model-quality} là \textbf{khoảng cách giữa
các mô hình rất hẹp}. Chênh lệch $\AUC$ giữa mô hình được chọn triển khai
(\BestModel{}) và hồi quy logistic chỉ là \GapBestLR{} -- chưa tới nửa điểm phần
trăm. Trong khi đó, hồi quy logistic có ba ưu thế rõ rệt: thời gian huấn luyện
ngắn hơn hàng chục lần, mỗi hệ số có thể diễn giải trực tiếp, và xác suất đầu ra
vốn đã được hiệu chuẩn tốt mà không cần bước hiệu chuẩn bổ sung.

Kết quả này phù hợp với phát hiện của Lessmann và cộng sự~\cite{lessmann2015}
trên các bộ dữ liệu tín dụng thực: các mô hình học tập hợp cho kết quả tốt hơn
nhưng mức chênh lệch so với hồi quy logistic không lớn.

\subsubsection{Giới hạn độ phân giải của phép đo}
\label{subsec:do-phan-giai}

Khoảng cách hẹp giữa các mô hình đặt ra một câu hỏi phương pháp mà phần lớn
nghiên cứu so sánh thuật toán không nêu: \textit{chênh lệch $\AUC$ nhỏ tới đâu
thì còn mang thông tin?}

Để trả lời, đề án chạy lại toàn bộ bước huấn luyện trên cùng tập dữ liệu, cùng
siêu tham số và cùng hạt giống ngẫu nhiên. Kết quả tái lập chính xác tới bốn chữ
số thập phân ở cả \NModels{} mô hình. Tuy nhiên khi chạy lại \textit{bước tinh
chỉnh siêu tham số}, giá trị $\AUC$ kiểm định chéo đo được dao động khoảng
$3 \times 10^{-4}$ giữa các lần chạy. Nguyên nhân mang tính kỹ thuật thuần túy:
các thư viện boosting dựng biểu đồ tần suất song song nhiều luồng, nên thứ tự
cộng dồn số thực dấu phẩy động thay đổi theo cách hệ điều hành cấp phát luồng.

Mức dao động này nhỏ, nhưng không nhỏ hơn khoảng cách giữa các mô hình dẫn đầu.
Hệ quả trực tiếp đã được quan sát trong quá trình thực nghiệm: hai bộ siêu tham
số của XGBoost có $\AUC$ kiểm định chéo cách nhau $1{,}6 \times 10^{-4}$ đã đổi
chỗ cho nhau giữa hai lần chạy, kéo theo thứ hạng của chính mô hình đó trên tập
kiểm thử cũng thay đổi.

Kết luận rút ra có ý nghĩa với mọi nghiên cứu so sánh thuật toán trên bài toán
này: \textbf{thứ hạng ở chữ số thập phân thứ tư của $\AUC$ không phải là bằng
chứng về sự vượt trội}. Các mô hình chênh nhau dưới ngưỡng đó cần được coi là
tương đương, và việc lựa chọn phải dựa trên các tiêu chí khác -- chi phí tính
toán, khả năng giải thích, độ ổn định -- chứ không dựa trên bảng xếp hạng.

\subsubsection{Đánh đổi giữa chất lượng và chi phí tính toán}

Bảng~\ref{tab:model-cost} trình bày chi phí huấn luyện và thời gian chấm điểm một
hồ sơ của từng mô hình.

\begin{table}[!htbp]
\centering
\caption{Chi phí huấn luyện và thời gian chấm điểm một hồ sơ}
\label{tab:model-cost}
\small
\setlength{\tabcolsep}{3.0pt}
\begin{tabular}{@{}L{3.799cm}L{2.849cm}R{1.899cm}R{1.994cm}R{1.994cm}R{1.709cm}@{}}
\toprule
\textbf{Mô hình} & \textbf{Nhóm} & \textbf{Thời gian HL (s)} & \textbf{Chấm điểm p50 (ms)} & \textbf{Chấm điểm p95 (ms)} & \textbf{Số vòng lặp} \\
\midrule
Kết hợp xếp tầng (Stacking) & Kết hợp xếp tầng & 214,8 & 141,38 & 202,39 & -- \\
CatBoost & Học tập hợp tuần tự & 91,5 & 4,61 & 10,46 & 353 \\
XGBoost & Học tập hợp tuần tự & 20,7 & 16,78 & 18,53 & 706 \\
LightGBM & Học tập hợp tuần tự & 10,8 & 7,85 & 8,81 & 156 \\
Hồi quy logistic & Tuyến tính & 3,2 & 10,63 & 16,55 & 94 \\
Gradient Boosting (histogram) & Học tập hợp tuần tự & 50,5 & 13,14 & 23,23 & 184 \\
Thẻ điểm WOE + Hồi quy logistic & Tuyến tính & 7,6 & 11,12 & 22,77 & 74 \\
Mạng nơ-ron truyền thẳng (MLP) & Mạng nơ-ron & 24,1 & 6,42 & 13,44 & 35 \\
Máy vectơ hỗ trợ (SVM-RBF) & Máy vectơ hỗ trợ & 16,1 & 15,10 & 20,64 & -- \\
Rừng ngẫu nhiên & Học tập hợp song song & 24 & 188,81 & 236,25 & -- \\
Cây cực ngẫu nhiên (Extra Trees) & Học tập hợp song song & 11,8 & 383,00 & 516,63 & -- \\
Cây quyết định (CART) & Cây quyết định & 5,5 & 12,06 & 14,93 & -- \\
\bottomrule
\end{tabular}
\end{table}


Hình~\ref{fig:auc-vs-latency} biểu diễn quan hệ giữa chất lượng dự báo và thời
gian chấm điểm.

\begin{figure}[!htbp]
\centering
\cffig[0.92\textwidth]{auc_vs_latency}
\caption{Đánh đổi giữa chất lượng dự báo và thời gian chấm điểm một hồ sơ}
\label{fig:auc-vs-latency}
\end{figure}

Đồ thị cho thấy một số mô hình nằm rõ rệt phía dưới đường biên hiệu quả. Rõ nhất
là mô hình kết hợp xếp tầng: nó đổi \GapTopDeploy{} điểm $\AUC$ lấy thời gian
chấm điểm cao gấp \LatRatioTopDeploy{} lần (\TopAUCLatPFive{} ms so với
\BestLatPFive{} ms ở phân vị 95). Extra Trees và rừng ngẫu nhiên vừa chậm hơn vừa
kém chính xác hơn nhóm boosting. Với bài toán kiểm duyệt trực tuyến -- nơi kết
quả phải trả về trong lúc cán bộ tín dụng đang làm việc với hồ sơ -- những mô
hình này không phải lựa chọn hợp lý dù chỉ xét riêng góc độ kỹ thuật.

\subsubsection{Tiêu chí chọn mô hình triển khai}
\label{subsec:kq-chon-mo-hinh}

Hai phân tích trên dẫn tới một quy tắc lựa chọn cụ thể, thay cho cách lấy thẳng
mô hình đứng đầu bảng xếp hạng $\AUC$. Quy tắc gồm hai bước:

\begin{enumerate}[label=(\arabic*)]
  \item \textbf{Xác định nhóm tương đương về chất lượng:} mọi mô hình có $\AUC$
  nằm trong khoảng \TieTolerance{} so với mô hình dẫn đầu. Ngưỡng này được chọn
  lớn hơn một bậc so với dao động đo đạc đã xác định ở
  mục~\ref{subsec:do-phan-giai}, đồng thời vẫn nhỏ hơn mức chênh lệch có ý nghĩa
  nghiệp vụ.
  \item \textbf{Trong nhóm đó, chọn mô hình có thời gian chấm điểm thấp nhất} --
  tức mô hình nằm trên biên Pareto của Hình~\ref{fig:auc-vs-latency}.
\end{enumerate}

Áp dụng trên kết quả thực nghiệm, nhóm tương đương gồm \NTiedModels{} mô hình:
kết hợp xếp tầng (\TopAUC{}), CatBoost (\CatAUC{}) và XGBoost (\XGBAUC{}). Thời
gian chấm điểm ở phân vị 95 của ba mô hình lần lượt là \StackLat{} ms,
\CatLat{} ms và \XGBLat{} ms. Mô hình được chọn triển khai là \BestModel{}.

Lựa chọn này có hai căn cứ. Thứ nhất, nếu lấy mô hình kết hợp xếp tầng, thời gian
xử lý một hồ sơ sẽ vượt mục tiêu 100 ms đặt ra tại
mục~\ref{subsec:yeu-cau} -- nghĩa là đánh đổi một yêu cầu đã công bố lấy
\GapTopDeploy{} điểm $\AUC$ nằm trong vùng nhiễu đo. Thứ hai, quy tắc này cho kết
quả \textit{ổn định}: khi thứ hạng trong nhóm dẫn đầu đảo chiều giữa các lần chạy,
mô hình được chọn không đổi, vì tiêu chí quyết định là thời gian chấm điểm --
đại lượng chênh nhau hàng chục lần chứ không phải vài phần vạn.

\subsubsection{Tác động nghiệp vụ}

Bảng~\ref{tab:model-business} quy đổi chất lượng thống kê thành các chỉ tiêu
nghiệp vụ tại ngưỡng tối ưu chi phí.

\begin{table}[!htbp]
\centering
\caption{Tác động nghiệp vụ của các mô hình tại ngưỡng tối ưu chi phí}
\label{tab:model-business}
\small
\setlength{\tabcolsep}{3.0pt}
\begin{tabular}{@{}L{2.718cm}R{1.500cm}R{1.521cm}R{1.565cm}R{1.500cm}R{1.750cm}R{1.750cm}R{1.521cm}@{}}
\toprule
\textbf{Mô hình} & \textbf{Ngưỡng} & \textbf{Tỷ lệ duyệt} & \textbf{Nợ xấu sau duyệt} & \textbf{Chi phí} & \textbf{Accuracy} & \textbf{Precision} & \textbf{Recall} \\
\midrule
\textbf{Kết hợp xếp tầng (Stacking)} & \textbf{0,0501} & \textbf{44,0\%} & \textbf{2,98\%} & \textbf{0,5532} & \textbf{0,5910} & \textbf{0,2925} & \textbf{0,9258} \\
CatBoost & 0,0746 & 50,5\% & 3,52\% & 0,5497 & 0,6459 & 0,3210 & 0,8994 \\
XGBoost & 0,0789 & 49,2\% & 3,32\% & 0,5431 & 0,6366 & 0,3161 & 0,9076 \\
LightGBM & 0,0758 & 45,2\% & 3,19\% & 0,5586 & 0,6003 & 0,2965 & 0,9183 \\
Hồi quy logistic & 0,0729 & 43,4\% & 2,94\% & 0,5550 & 0,5856 & 0,2899 & 0,9277 \\
Gradient Boosting (histogram) & 0,0749 & 44,7\% & 3,21\% & 0,5624 & 0,5952 & 0,2938 & 0,9189 \\
Thẻ điểm WOE + Hồi quy logistic & 0,0848 & 48,0\% & 3,26\% & 0,5464 & 0,6259 & 0,3101 & 0,9114 \\
Mạng nơ-ron truyền thẳng (MLP) & 0,0778 & 48,7\% & 3,40\% & 0,5510 & 0,6311 & 0,3126 & 0,9063 \\
Máy vectơ hỗ trợ (SVM-RBF) & 0,0702 & 46,1\% & 3,28\% & 0,5588 & 0,6074 & 0,2999 & 0,9145 \\
Rừng ngẫu nhiên & 0,0682 & 49,4\% & 3,62\% & 0,5618 & 0,6350 & 0,3140 & 0,8988 \\
Cây cực ngẫu nhiên (Extra Trees) & 0,0866 & 42,8\% & 3,17\% & 0,5714 & 0,5777 & 0,2854 & 0,9233 \\
Cây quyết định (CART) & 0,0707 & 38,1\% & 4,08\% & 0,6444 & 0,5267 & 0,2605 & 0,9120 \\
\bottomrule
\end{tabular}

\vspace{2pt}
{\footnotesize\itshape Ngưỡng được chọn theo hàm chi phí với tỷ lệ $c_{FN}:c_{FP} = 12:1$ trên tập kiểm định.}
\end{table}


\subsubsection{Mức độ quá khớp và độ ổn định theo thời gian}
\label{subsec:kq-on-dinh}

\paragraph{Mức độ quá khớp.} Bảng~\ref{tab:overfit} đối chiếu $\AUC$ trên tập
kiểm định với $\AUC$ trên tập kiểm thử. Chênh lệch trung bình giữa hai tập là
\MeanOverfitGap{} và mức lớn nhất là \MaxOverfitGap{} -- đều rất nhỏ so với bản
thân giá trị $\AUC$. Kết quả này cho thấy bốn cơ chế kiểm soát quá khớp trình bày
ở mục~\ref{subsec:kiem-soat-qua-khop} đã phát huy tác dụng: không mô hình nào phụ
thuộc quá mức vào đặc thù của giai đoạn huấn luyện.

\begin{table}[!htbp]
\centering
\caption{Mức độ quá khớp đo bằng chênh lệch AUC giữa hai tập dữ liệu}
\label{tab:overfit}
\small
\setlength{\tabcolsep}{5.0pt}
\begin{tabular}{@{}L{5.4cm}R{2.9cm}R{2.9cm}R{2.6cm}@{}}
\toprule
\textbf{Mô hình} & \textbf{AUC (kiểm định)} & \textbf{AUC (kiểm thử)} & \textbf{Chênh lệch} \\
\midrule
Kết hợp xếp tầng (Stacking) & 0,8639 & 0,8620 & 0,0019 \\
CatBoost & 0,8633 & 0,8614 & 0,0019 \\
XGBoost & 0,8646 & 0,8612 & 0,0034 \\
LightGBM & 0,8595 & 0,8590 & 0,0005 \\
Hồi quy logistic & 0,8590 & 0,8579 & 0,0011 \\
Gradient Boosting (histogram) & 0,8616 & 0,8575 & 0,0041 \\
Thẻ điểm WOE + Hồi quy logistic & 0,8610 & 0,8570 & 0,0040 \\
Mạng nơ-ron truyền thẳng (MLP) & 0,8554 & 0,8560 & -0,0006 \\
Máy vectơ hỗ trợ (SVM-RBF) & 0,8578 & 0,8540 & 0,0038 \\
Rừng ngẫu nhiên & 0,8565 & 0,8527 & 0,0038 \\
Cây cực ngẫu nhiên (Extra Trees) & 0,8521 & 0,8471 & 0,0050 \\
Cây quyết định (CART) & 0,8231 & 0,8206 & 0,0025 \\
\bottomrule
\end{tabular}

\vspace{2pt}
{\footnotesize\itshape Chênh lệch dương lớn cho thấy mô hình phụ thuộc quá mức vào đặc thù của giai đoạn huấn luyện.}
\end{table}


\paragraph{Trôi dạt phân phối dữ liệu.} Bảng~\ref{tab:drift} trình bày mười đặc
trưng có mức trôi dạt lớn nhất. Toàn bộ tập đặc trưng có $\PSI$ dưới ngưỡng cảnh
báo, với giá trị lớn nhất là \MaxPSI{}; số đặc trưng ở mức cảnh báo là
\NPSIWarn{} và ở mức báo động là \NPSIAlert{}.

\begin{table}[!htbp]
\centering
\caption{Mười đặc trưng có mức trôi dạt phân phối lớn nhất}
\label{tab:drift}
\small
\setlength{\tabcolsep}{5.0pt}
\begin{tabular}{@{}L{6.4cm}R{2.6cm}L{4.4cm}@{}}
\toprule
\textbf{Đặc trưng} & \textbf{PSI} & \textbf{Trạng thái} \\
\midrule
salary\_\allowbreak cnt\_\allowbreak 90d & 0,0598 & ỔN ĐỊNH \\
salary\_\allowbreak cnt\_\allowbreak 30d & 0,0571 & ỔN ĐỊNH \\
income\_\allowbreak stability & 0,0533 & ỔN ĐỊNH \\
salary\_\allowbreak cnt\_\allowbreak 180d & 0,0504 & ỔN ĐỊNH \\
income\_\allowbreak verification\_\allowbreak ratio & 0,0395 & ỔN ĐỊNH \\
outflow\_\allowbreak ratio\_\allowbreak 30d & 0,0388 & ỔN ĐỊNH \\
salary\_\allowbreak amt\_\allowbreak 30d & 0,0374 & ỔN ĐỊNH \\
net\_\allowbreak flow\_\allowbreak 30d & 0,0363 & ỔN ĐỊNH \\
txn\_\allowbreak amt\_\allowbreak std\_\allowbreak 30d & 0,0336 & ỔN ĐỊNH \\
inflow\_\allowbreak to\_\allowbreak installment & 0,0291 & ỔN ĐỊNH \\
\bottomrule
\end{tabular}

\vspace{2pt}
{\footnotesize\itshape Ngưỡng cảnh báo PSI $= 0{,}10$; ngưỡng báo động PSI $= 0{,}25$.}
\end{table}


\paragraph{Chất lượng dự báo theo thời gian.} Bảng~\ref{tab:perf-time} theo dõi
chất lượng dự báo theo từng tháng của giai đoạn kiểm thử. Giá trị $\AUC$ dao động
trong khoảng hẹp quanh mức 0,86 mà không có xu hướng giảm rõ rệt, và chênh lệch
giữa tỷ lệ nợ xấu thực tế với mức $\PD$ trung bình dự báo luôn dưới một điểm phần
trăm. Trong phạm vi bốn tháng quan sát, mô hình giữ được cả khả năng phân biệt
lẫn mức hiệu chuẩn. Tuy nhiên, bốn tháng là khoảng thời gian ngắn so với chu kỳ
suy giảm điển hình của mô hình tín dụng trong thực tế, nên kết luận về độ ổn định
chỉ có giá trị trong phạm vi dữ liệu thực nghiệm.

\begin{table}[!htbp]
\centering
\caption{Chất lượng dự báo của mô hình theo từng tháng của giai đoạn kiểm thử}
\label{tab:perf-time}
\small
\setlength{\tabcolsep}{3.0pt}
\begin{tabular}{@{}L{1.671cm}R{1.495cm}R{2.023cm}R{2.287cm}R{1.847cm}R{1.500cm}R{1.500cm}R{1.500cm}@{}}
\toprule
\textbf{Tháng} & \textbf{Số hồ sơ} & \textbf{Nợ xấu thực tế (\%)} & \textbf{PD dự báo trung bình (\%)} & \textbf{Chênh lệch (đpt)} & \textbf{AUC} & \textbf{KS} & \textbf{Gini} \\
\midrule
2026-03 & 1.652 & 17,86 & 17,76 & -0,09 & 0,8665 & 0,5953 & 0,7330 \\
2026-04 & 2.485 & 17,46 & 17,13 & -0,33 & 0,8584 & 0,5730 & 0,7168 \\
2026-05 & 2.429 & 18,53 & 17,74 & -0,79 & 0,8551 & 0,5673 & 0,7102 \\
2026-06 & 2.434 & 16,93 & 17,26 & 0,33 & 0,8679 & 0,5782 & 0,7358 \\
\bottomrule
\end{tabular}

\vspace{2pt}
{\footnotesize\itshape Chênh lệch tính bằng điểm phần trăm giữa mức PD dự báo trung bình và tỷ lệ nợ xấu thực tế của tháng.}
\end{table}


\begin{figure}[!htbp]
\centering
\cffig[0.95\textwidth]{performance_over_time}
\caption{Diễn biến chất lượng dự báo theo tháng trên giai đoạn kiểm thử}
\label{fig:perf-time}
\end{figure}

\section{Mô phỏng và giải thích mô hình}
\label{sec:mo-phong-he-thong}

\subsection{Thử nghiệm giải thích quyết định mô hình}
\label{subsec:kq-giai-thich}

\subsubsection{Giải thích ở cấp độ toàn cục}

Phân tích SHAP được thực hiện trên mẫu đại diện của tập kiểm thử.
Bảng~\ref{tab:shap-group} tổng hợp đóng góp theo nhóm nguồn dữ liệu.

\begin{table}[!htbp]
\centering
\caption{Tỷ trọng đóng góp vào quyết định của mô hình theo nhóm đặc trưng}
\label{tab:shap-group}
\small
\setlength{\tabcolsep}{3.0pt}
\begin{tabular}{@{}L{3.648cm}R{1.911cm}R{1.911cm}R{2.258cm}R{2.085cm}R{2.432cm}@{}}
\toprule
\textbf{Nhóm đặc trưng} & \textbf{Số đặc trưng} & \textbf{Tổng |SHAP|} & \textbf{|SHAP| trung bình} & \textbf{|SHAP| lớn nhất} & \textbf{Tỷ trọng đóng góp (\%)} \\
\midrule
Hành vi giao dịch & 76 & 1,2895 & 0,0170 & 0,1623 & 34,00 \\
Lịch sử tín dụng (CIC) & 14 & 1,1692 & 0,0835 & 0,4217 & 30,83 \\
Biến tổ hợp & 12 & 0,9091 & 0,0758 & 0,1899 & 23,97 \\
Thông tin cá nhân & 8 & 0,2662 & 0,0333 & 0,0909 & 7,02 \\
Tham số khoản vay & 6 & 0,1587 & 0,0264 & 0,0432 & 4,18 \\
\bottomrule
\end{tabular}

\vspace{2pt}
{\footnotesize\itshape Tính bằng tổng giá trị SHAP tuyệt đối trung bình của các đặc trưng trong từng nhóm, chuẩn hóa về tổng bằng 100\%.}
\end{table}


Kết quả cho thấy nhóm đặc trưng hành vi giao dịch đóng góp
\ShapBehaviourPct{}\% vào quyết định của mô hình -- cao hơn nhóm CIC
(\ShapCICPct{}\%) vốn được coi là nguồn thông tin truyền thống quan trọng nhất.
Nhóm biến tổ hợp đóng góp \ShapDerivedPct{}\%, còn hai nhóm thông tin cá nhân và
tham số khoản vay chỉ đóng góp lần lượt \ShapStaticPct{}\% và \ShapLoanPct{}\%.

Phát hiện này có ý nghĩa thực tiễn: dữ liệu giao dịch là nguồn mà tổ chức tín
dụng tự sở hữu, không phát sinh chi phí truy vấn bên ngoài và luôn có sẵn kể cả
với nhóm khách hàng thiếu lịch sử tín dụng. Việc nhóm này đóng góp phần lớn nhất
vào quyết định của mô hình là căn cứ định lượng cho khuyến nghị khai thác dữ liệu
giao dịch trong công tác kiểm duyệt tín dụng.

Kết quả trên cần được đặt cạnh phân tích $\IV$ ở mục~\ref{subsec:kq-dac-trung} để
tránh diễn giải sai. Xét từng đặc trưng riêng lẻ, biến CIC mạnh hơn hẳn biến giao
dịch. Nhưng trong mô hình hoàn chỉnh, tổng đóng góp của nhóm giao dịch lại lớn
hơn, vì nhóm này gồm nhiều đặc trưng bổ sung thông tin cho nhau ở những khía cạnh
khác nhau của hành vi tài chính.

\begin{figure}[!htbp]
\centering
\cffig[0.95\textwidth]{shap_summary}
\caption{Biểu đồ tóm tắt SHAP của hai mươi đặc trưng quan trọng nhất}
\label{fig:shap-summary}
\end{figure}

\begin{figure}[!htbp]
\centering
\cffig[0.9\textwidth]{shap_group_share}
\caption{Tỷ trọng đóng góp vào quyết định theo nhóm nguồn dữ liệu}
\label{fig:shap-group}
\end{figure}

\subsubsection{Kiểm chứng chéo bằng phép hoán vị}

Kết quả SHAP được kiểm chứng độc lập bằng phép hoán vị: với mỗi đặc trưng, giá
trị của nó được xáo trộn ngẫu nhiên và mức suy giảm $\AUC$ được ghi nhận. Thứ
hạng các đặc trưng theo hai phương pháp về cơ bản trùng khớp ở nhóm đầu bảng, xác
nhận rằng kết quả SHAP không phải là sản phẩm của một đặc thù kỹ thuật nào đó
trong cách tính.

\begin{figure}[!htbp]
\centering
\cffig[0.95\textwidth]{permutation_importance_bar}
\caption{Tầm quan trọng đặc trưng theo phép hoán vị}
\label{fig:permutation}
\end{figure}

\subsubsection{Giải thích ở cấp độ từng hồ sơ}

Với mỗi hồ sơ, mô hình sinh danh sách lý do gồm ba thành phần: luật nghiệp vụ bị
vi phạm, chỉ tiêu định lượng then chốt kèm giá trị thực tế, và các đặc trưng có
đóng góp SHAP lớn nhất. Danh sách này được trả về cùng kết quả chấm điểm, cho
phép cán bộ tín dụng đánh giá tính hợp lý của khuyến nghị thay vì phải chấp nhận
một con số không có căn cứ.

\subsection{Mô phỏng xác định ngưỡng phê duyệt tín dụng}
\label{subsec:kq-nguong}

\subsubsection{So sánh ba tiêu chí xác định điểm cắt}

Bảng~\ref{tab:cutoff} trình bày kết quả của ba tiêu chí xác định điểm cắt, tính
trên tập kiểm định.

\begin{table}[!htbp]
\centering
\caption{So sánh ba tiêu chí xác định điểm cắt phê duyệt}
\label{tab:cutoff}
\small
\setlength{\tabcolsep}{3.0pt}
\begin{tabular}{@{}L{4.255cm}R{1.850cm}R{1.850cm}R{2.220cm}R{2.220cm}R{1.850cm}@{}}
\toprule
\textbf{Tiêu chí} & \textbf{Điểm cắt PD} & \textbf{Tỷ lệ duyệt (\%)} & \textbf{Nợ xấu nhóm duyệt (\%)} & \textbf{Nợ xấu nhóm từ chối (\%)} & \textbf{Chi phí kỳ vọng} \\
\midrule
Tối thiểu chi phí (c\_\allowbreak FN:c\_\allowbreak FP = 12:1) & 0,07 & 51,64 & 3,33 & 31,23 & 0,54 \\
Tỷ lệ duyệt mục tiêu 70\% & 0,17 & 72,30 & 6,01 & 45,05 & 0,67 \\
Trần nợ xấu nhóm duyệt 5\% & 0,15 & 65,30 & 4,92 & 39,23 & 0,60 \\
\bottomrule
\end{tabular}
\end{table}


Ba tiêu chí cho ba điểm cắt khác nhau rõ rệt, phản ánh ba cách đặt bài toán khác
nhau. Tiêu chí tối thiểu chi phí cho điểm cắt thấp nhất (\CostOptimalCut{}) và do
đó thận trọng nhất, với tỷ lệ nợ xấu trong nhóm được duyệt chỉ
\CutoffCostBadRate{}. Tiêu chí tỷ lệ duyệt mục tiêu 70\% cho điểm cắt cao nhất,
đổi lại tỷ lệ nợ xấu nhóm duyệt tăng lên \CutoffRateBadRate{}. Tiêu chí trần nợ
xấu 5\% nằm ở giữa, với tỷ lệ nợ xấu nhóm duyệt \CutoffCapBadRate{}.

Việc trình bày song song ba tiêu chí -- thay vì áp đặt một ngưỡng duy nhất -- là
có chủ đích: lựa chọn giữa chúng là \textit{quyết định chính sách} thuộc thẩm
quyền của tổ chức tín dụng, không phải quyết định kỹ thuật của người xây dựng mô
hình.

\subsubsection{Độ nhạy của kết quả theo điểm cắt}

Bảng~\ref{tab:cutoff-sens} và Hình~\ref{fig:cutoff-sens} định lượng đánh đổi giữa
quy mô cấp tín dụng và chất lượng danh mục.

\begin{table}[!htbp]
\centering
\caption{Độ nhạy của kết quả kiểm duyệt theo điểm cắt}
\label{tab:cutoff-sens}
\small
\setlength{\tabcolsep}{3.0pt}
\begin{tabular}{@{}R{2.0cm}R{2.2cm}R{2.2cm}R{2.6cm}R{2.4cm}R{2.2cm}@{}}
\toprule
\textbf{Điểm cắt PD} & \textbf{Số hồ sơ duyệt} & \textbf{Tỷ lệ duyệt (\%)} & \textbf{Nợ xấu nhóm duyệt (\%)} & \textbf{Nợ xấu bị chặn (\%)} & \textbf{Chi phí kỳ vọng} \\
\midrule
0.02 & 2.085 & 23,17 & 1,10 & 21,56 & 0,63 \\
0.04 & 2.775 & 30,83 & 1,66 & 23,58 & 0,59 \\
0.06 & 3.505 & 38,94 & 2,25 & 26,11 & 0,56 \\
0.08 & 4.893 & 54,37 & 3,56 & 32,63 & 0,54 \\
0.1 & 4.893 & 54,37 & 3,56 & 32,63 & 0,54 \\
0.15 & 5.877 & 65,30 & 4,92 & 39,23 & 0,60 \\
0.2 & 6.843 & 76,03 & 6,62 & 49,19 & 0,73 \\
0.25 & 6.999 & 77,77 & 6,94 & 51,37 & 0,76 \\
0.32 & 7.533 & 83,70 & 8,56 & 59,24 & 0,93 \\
0.4 & 7.623 & 84,70 & 8,87 & 60,86 & 0,96 \\
\bottomrule
\end{tabular}
\end{table}


\begin{figure}[!htbp]
\centering
\cffig[0.9\textwidth]{cutoff_sensitivity}
\caption{Đánh đổi giữa tỷ lệ duyệt và tỷ lệ nợ xấu theo điểm cắt}
\label{fig:cutoff-sens}
\end{figure}

Đường cong này là công cụ đối thoại giữa bộ phận mô hình và bộ phận kinh doanh:
thay vì tranh luận về một giá trị ngưỡng trừu tượng, mỗi lựa chọn được quy đổi
thành hệ quả cụ thể về tỷ lệ duyệt và tỷ lệ nợ xấu.

\subsubsection{Hai ngưỡng cho ba vùng quyết định}

Từ điểm cắt tối ưu chi phí, hai ngưỡng phân vùng được xác định:
$\tau_{\text{duyệt}} = \PDApproveMax{}$ và
$\tau_{\text{từ chối}} = \PDRejectMin{}$. Kết quả áp dụng trên tập kiểm thử được
trình bày tại Bảng~\ref{tab:decision-mix}.

\begin{table}[!htbp]
\centering
\caption{Cơ cấu khuyến nghị của mô hình và tỷ lệ nợ xấu thực tế theo từng nhóm}
\label{tab:decision-mix}
\small
\setlength{\tabcolsep}{5.0pt}
\begin{tabular}{@{}L{5.854cm}R{2.537cm}R{2.537cm}R{3.317cm}@{}}
\toprule
\textbf{Kết quả} & \textbf{Số hồ sơ} & \textbf{Tỷ trọng (\%)} & \textbf{Tỷ lệ nợ xấu thực tế (\%)} \\
\midrule
Khuyến nghị duyệt & 1.459 & 29,18 & 3,43 \\
Chuyển thẩm định thủ công & 1.637 & 32,74 & 10,93 \\
Khuyến nghị từ chối & 1.904 & 38,08 & 34,45 \\
\bottomrule
\end{tabular}
\end{table}


Kết quả xác nhận mô hình hoạt động đúng thiết kế: tỷ lệ nợ xấu thực tế tăng đơn
điệu theo ba nhóm khuyến nghị, lần lượt là \BadApprove{} ở nhóm khuyến nghị
duyệt, \BadReview{} ở nhóm chuyển thẩm định và \BadReject{} ở nhóm khuyến nghị từ
chối. Mức chênh lệch hơn mười lần giữa nhóm đầu và nhóm cuối cho thấy khả năng
phân tách thực tế của mô hình.

Tỷ lệ hồ sơ được xử lý tự động đạt \AutomationRate{}, nghĩa là khoảng một phần ba
số hồ sơ (\ReviewRate{}) vẫn cần cán bộ tín dụng xem xét. Đây là kết quả phù hợp
với định hướng \textit{hỗ trợ} của đề án: những hồ sơ mà mô hình không đủ chắc
chắn được chuyển cho con người thay vì bị quyết định máy móc.

\begin{figure}[!htbp]
\centering
\cffig[0.98\textwidth]{decision_mix}
\caption{Cơ cấu khuyến nghị và tỷ lệ nợ xấu thực tế theo từng nhóm}
\label{fig:decision-mix}
\end{figure}

\subsubsection{Tần suất vi phạm luật nghiệp vụ}

Bảng~\ref{tab:rule-hits} trình bày tần suất vi phạm của \NRules{} luật nghiệp vụ
trong quá trình mô phỏng.

\begin{table}[!htbp]
\centering
\caption{Tần suất vi phạm các luật nghiệp vụ}
\label{tab:rule-hits}
\small
\setlength{\tabcolsep}{5.0pt}
\begin{tabular}{@{}L{2.8cm}L{6.0cm}L{2.6cm}R{2.6cm}@{}}
\toprule
\textbf{Mã luật} & \textbf{Tên luật} & \textbf{Hành động} & \textbf{Tỷ lệ vi phạm} \\
\midrule
R501 & Dòng tiền vào thực tế thấp hơn nhiều so với thu nhập khai báo & REVIEW & 23,08\% \\
R202 & Thu nhập còn lại sau trả nợ không đủ chi phí sinh hoạt & HARD\_\allowbreak REJECT & 21,10\% \\
R104 & Số ngày quá hạn tối đa 12 tháng vượt 30 ngày & REVIEW & 12,98\% \\
R101 & Nhóm nợ xấu nhất 12 tháng vượt nhóm 2 & HARD\_\allowbreak REJECT & 11,74\% \\
R107 & Tỷ lệ sử dụng hạn mức trên 95\% & WARNING & 11,10\% \\
R103 & Có nợ đã xử lý rủi ro / xóa nợ & HARD\_\allowbreak REJECT & 10,18\% \\
R401 & Không lấy được dữ liệu CIC & REVIEW & 9,74\% \\
R402 & Khách hàng thiếu lịch sử tín dụng (thin-file) & WARNING & 9,74\% \\
R002 & Thu nhập dưới mức tối thiểu 6 triệu VND & HARD\_\allowbreak REJECT & 9,28\% \\
R102 & Đang có nợ nhóm 3 trở lên & HARD\_\allowbreak REJECT & 9,28\% \\
R201 & Tỷ lệ nợ trên thu nhập (DTI) vượt 60\% & HARD\_\allowbreak REJECT & 7,46\% \\
R502 & Tỷ trọng chi tiêu nhóm rủi ro cao bất thường & REVIEW & 6,04\% \\
\bottomrule
\end{tabular}
\end{table}


\begin{figure}[!htbp]
\centering
\cffig[0.95\textwidth]{rule_hits}
\caption{Tần suất vi phạm các luật nghiệp vụ}
\label{fig:rule-hits}
\end{figure}

\subsection{Thử nghiệm tích hợp API và luồng xử lý hồ sơ}
\label{subsec:kq-api}

\subsubsection{Kết quả tích hợp dịch vụ REST}

Hệ thống được đóng gói thành một dịch vụ REST với năm điểm cuối. Kết quả thử
nghiệm tích hợp trình bày tại Bảng~\ref{tab:api}.

\begin{table}[!htbp]
\centering
\caption{Kết quả thử nghiệm tích hợp dịch vụ REST}
\label{tab:api}
\small
\setlength{\tabcolsep}{5.0pt}
\begin{tabular}{@{}L{4.600cm}L{4.161cm}R{1.702cm}L{3.783cm}@{}}
\toprule
\textbf{Điểm cuối} & \textbf{Mô tả} & \textbf{Mã HTTP} & \textbf{Kết quả} \\
\midrule
GET /health & Kiểm tra trạng thái dịch vụ & 200 & Đạt \\
POST /api/v1/screen & Kiểm duyệt từng hồ sơ (50 lần gọi) & 200 & 50/50 thành công, 62.4 ms/hồ sơ \\
POST /api/v1/screen/batch & Kiểm duyệt theo lô 50 hồ sơ & 200 & 50 hồ sơ, 2245 ms (44.79 ms/hồ sơ) \\
GET /api/v1/features/\{id\} & Truy xuất vector đặc trưng đầu vào & 200 & 157 đặc trưng \\
GET /admin/rules & Danh mục luật nghiệp vụ đang áp dụng & 200 & 21 luật \\
\bottomrule
\end{tabular}
\end{table}


Toàn bộ điểm cuối hoạt động đúng. Điểm cuối kiểm duyệt theo lô cho thời gian xử
lý bình quân mỗi hồ sơ thấp hơn so với gọi từng hồ sơ riêng lẻ, phản ánh phần chi
phí cố định của mỗi lần gọi dịch vụ được phân bổ trên nhiều hồ sơ.

\subsubsection{Thời gian xử lý một hồ sơ}

Bảng~\ref{tab:processing-time} phân rã thời gian xử lý một hồ sơ theo từng chặng.

\begin{table}[!htbp]
\centering
\caption{Thời gian xử lý một hồ sơ theo từng chặng}
\label{tab:processing-time}
\small
\setlength{\tabcolsep}{5.0pt}
\begin{tabular}{@{}L{4.6cm}R{2.6cm}R{2.2cm}R{2.2cm}R{2.2cm}@{}}
\toprule
\textbf{Chặng xử lý} & \textbf{Trung bình (ms)} & \textbf{p50 (ms)} & \textbf{p95 (ms)} & \textbf{p99 (ms)} \\
\midrule
Chấm điểm mô hình & 40,174 & 34,773 & 74,507 & 98,067 \\
Áp dụng luật nghiệp vụ & 0,048 & 0,046 & 0,086 & 0,126 \\
Áp dụng chính sách & 0,053 & 0,051 & 0,088 & 0,154 \\
Tổng cộng & 40,327 & 34,928 & 74,657 & 98,372 \\
\bottomrule
\end{tabular}

\vspace{2pt}
{\footnotesize\itshape Đo trên một tiến trình đơn, không tính thời gian truyền qua mạng và thời gian truy vấn dữ liệu từ hệ thống nguồn.}
\end{table}


Kết quả cho thấy \textbf{chấm điểm mô hình chiếm gần như toàn bộ} thời gian xử
lý: trung bình \SysMean{} ms cho một hồ sơ, trong đó phần áp dụng luật nghiệp vụ
và chính sách chỉ chiếm chưa tới 0,1 ms. Phân vị 95\% của tổng thời gian là
\SysPNinetyFive{} ms và phân vị 99\% là \SysPNinetyNine{} ms, đều nằm dưới mục
tiêu 100 ms đã đặt ra ở mục~\ref{subsec:yeu-cau}.

Thời gian chấm điểm trong mô phỏng quy trình (\SysMean{} ms) cao hơn số đo ở bước
huấn luyện (\BestLatPFive{} ms tại phân vị 95\%) do hai phép đo thực hiện trong
điều kiện khác nhau: bước huấn luyện đo trên ma trận đặc trưng đã chuẩn bị sẵn,
còn mô phỏng quy trình đo toàn bộ đường xử lý từ dữ liệu đầu vào thô, bao gồm cả
việc dựng khung dữ liệu, ép kiểu và ánh xạ hạng mục. Phần chênh lệch chính là chi
phí của lớp tích hợp. Đây là kết quả có giá trị thực tiễn: khi triển khai, chi
phí chuẩn bị dữ liệu đầu vào có thể lớn hơn chi phí suy luận của bản thân mô
hình, nên việc tối ưu hiệu năng cần bắt đầu từ khâu này.

\section{Đánh giá chung}
\label{sec:danh-gia-chung}

\subsection{So sánh với quy trình thẩm định truyền thống}
\label{subsec:so-sanh-truyen-thong}

Để đánh giá giá trị thực tế của mô hình, đề án mô phỏng quy trình thẩm định
truyền thống bằng chính bộ luật nghiệp vụ đang áp dụng, không sử dụng mô hình học
máy: hồ sơ vi phạm luật loại trừ cứng bị từ chối, hồ sơ vi phạm luật cảnh báo
được chuyển thẩm định, phần còn lại được duyệt. Hai phương pháp được chạy trên
cùng một tập hồ sơ; kết quả trình bày tại Bảng~\ref{tab:traditional}.

\begin{table}[!htbp]
\centering
\caption{So sánh mô hình đề xuất với quy trình thẩm định truyền thống}
\label{tab:traditional}
\small
\setlength{\tabcolsep}{5.0pt}
\begin{tabular}{@{}L{6.312cm}R{4.142cm}R{4.142cm}@{}}
\toprule
\textbf{Chỉ tiêu} & \textbf{Truyền thống} & \textbf{Đề xuất} \\
\midrule
Tỷ lệ duyệt (\%) & 35.82 & 29.18 \\
Tỷ lệ thẩm định thủ công (\%) & 27.22 & 32.74 \\
Tỷ lệ từ chối (\%) & 36.96 & 38.08 \\
Nợ xấu nhóm duyệt (\%) & 6.14 & 3.43 \\
Nợ xấu chặn được (\%) & 33.93 & 34.45 \\
Tỷ lệ tự động hóa (\%) & 72.78 & 67.26 \\
Thời gian xử lý (ms/hồ sơ) & 0.262 & 40.327 \\
\bottomrule
\end{tabular}

\vspace{2pt}
{\footnotesize\itshape Quy trình truyền thống được mô phỏng bằng bộ luật nghiệp vụ đang áp dụng, không sử dụng mô hình học máy.}
\end{table}


\begin{figure}[!htbp]
\centering
\cffig[0.95\textwidth]{traditional_comparison}
\caption{So sánh mô hình đề xuất với quy trình thẩm định truyền thống}
\label{fig:traditional}
\end{figure}

Kết quả phản ánh một đánh đổi gồm bốn khía cạnh.

\paragraph{Về chất lượng danh mục -- cải thiện rõ rệt.} Tỷ lệ nợ xấu trong nhóm
hồ sơ được duyệt giảm từ \TradBadRate{} xuống \PropBadRate{}, tức giảm
\BadRateReduction{} điểm phần trăm, tương đương mức giảm tương đối
\BadRateReductionPct{}. Đây là chỉ tiêu có ý nghĩa kinh tế trực tiếp: với cùng
quy mô dư nợ được cấp, tổn thất tín dụng kỳ vọng giảm gần một nửa.

\paragraph{Về quy mô cấp tín dụng -- giảm.} Đổi lại, tỷ lệ hồ sơ được duyệt giảm
từ \TradApproveRate{} xuống \ApproveRate{}. Hệ thống thận trọng hơn quy trình
truyền thống: nó phát hiện và loại ra những hồ sơ mà bộ quy tắc cố định bỏ sót.

\paragraph{Về khối lượng công việc thủ công -- tăng nhẹ.} Tỷ lệ hồ sơ chuyển thẩm
định thủ công tăng từ \TradReviewRate{} lên \PropReviewRate{}. Đây là hệ quả trực
tiếp của thiết kế ba vùng: mô hình chủ động thừa nhận vùng không chắc chắn của
mình thay vì đưa ra khuyến nghị dứt khoát cho mọi hồ sơ.

\paragraph{Về thời gian xử lý.} Quy trình chỉ dùng luật nhanh hơn nhiều về mặt
tính toán (dưới 0,2 ms so với \SysMean{} ms). Tuy nhiên so sánh này không phản
ánh thực tế vận hành: trong quy trình truyền thống, phần lớn thời gian nằm ở công
đoạn cán bộ tín dụng thu thập, đối chiếu và đánh giá hồ sơ -- tính bằng phút hoặc
giờ -- chứ không phải ở việc kiểm tra bộ quy tắc. Mức chênh lệch vài chục mili
giây giữa hai phương pháp là không đáng kể trên thang thời gian đó.

\paragraph{Diễn giải tổng hợp.} Đánh đổi mà mô hình mang lại là: chấp nhận giảm
\ApproveRateDrop{} điểm phần trăm tỷ lệ duyệt và tăng \ReviewRateRise{} điểm phần
trăm khối lượng thẩm định thủ công, để đổi lấy việc giảm gần một nửa tỷ lệ nợ xấu
trong nhóm được duyệt. Mức đánh đổi này có phù hợp hay không phụ thuộc khẩu vị
rủi ro và chiến lược kinh doanh của từng tổ chức. Khác biệt cơ bản so với bộ quy
tắc cố định là mô hình \textit{cho phép điều chỉnh} đánh đổi thông qua hai ngưỡng
phân vùng, trong khi bộ quy tắc không có cơ chế tương đương.

\subsection{Đối chiếu kết quả với các công trình nghiên cứu liên quan}
\label{subsec:doi-chieu-nghien-cuu}

\paragraph{Về thứ hạng giữa các thuật toán.} Kết quả của đề án phù hợp với kết
luận của Lessmann và cộng sự~\cite{lessmann2015}: nhóm học tập hợp đứng đầu nhưng
khoảng cách với hồi quy logistic không lớn. Chênh lệch $\AUC$ đo được trong đề án
là \GapBestLR{}, thậm chí còn hẹp hơn khoảng 0,01--0,03 mà công trình trên ghi
nhận trên các bộ dữ liệu tín dụng thực. Chênh lệch hẹp hơn có thể giải thích bằng
đặc điểm của dữ liệu mô phỏng: quan hệ giữa biến tiềm ẩn và nhãn được sinh theo
một hàm log-odds có cấu trúc tương đối trơn, nên phần quan hệ phi tuyến mà các mô
hình học tập hợp khai thác được so với mô hình tuyến tính là hạn chế. Trên dữ
liệu thực, nơi quan hệ giữa các biến phức tạp hơn, khoảng cách này nhiều khả năng
sẽ rộng hơn. Các tổng quan hệ thống gần đây~\cite{shi2022} cũng đi tới kết luận
tương tự về thứ hạng giữa các nhóm thuật toán.

\paragraph{Về vai trò của đặc trưng so với thuật toán.} Kết quả phân tích SHAP
của đề án củng cố nhận định phổ biến trong các tổng quan: yếu tố quyết định chất
lượng mô hình nằm ở chất lượng và độ phong phú của đặc trưng nhiều hơn ở việc lựa
chọn thuật toán. Việc nhóm đặc trưng hành vi giao dịch đóng góp \ShapBehaviourPct{}\%
-- vượt cả nhóm CIC truyền thống -- là minh chứng cụ thể cho nhận định này, đồng
thời phù hợp với quan sát của Onay và Öztürk~\cite{onay2018} về vai trò ngày càng
lớn của dữ liệu hành vi trong bối cảnh dữ liệu lớn.

\paragraph{Về khoảng cách so với các nghiên cứu trong nước.} Các nghiên cứu ứng
dụng học máy trong đánh giá rủi ro tín dụng tại Việt Nam~\cite{thu2025,thao2025}
chủ yếu tập trung vào việc so sánh chất lượng dự báo giữa các thuật toán trên dữ
liệu dạng bảng có sẵn. Đề án này bổ sung ba nội dung mà các công trình đó chưa đề
cập đầy đủ: (i) tích hợp mô hình vào một quy trình hỗ trợ kiểm duyệt hoàn chỉnh có
luật nghiệp vụ và cơ chế giải thích; (ii) đánh giá chất lượng hiệu chuẩn xác suất
-- điều kiện cần để dùng $\PD$ trong phân vùng quyết định; và (iii) đối chiếu
định lượng với quy trình thẩm định truyền thống.

\paragraph{Về điểm khác biệt trong thiết kế thực nghiệm.} Phần lớn nghiên cứu về
chấm điểm tín dụng đánh giá mô hình bằng cách chia dữ liệu ngẫu nhiên. Đề án sử
dụng kiểm định ngoài thời gian, cho kết quả thấp hơn nhưng sát thực tế hơn. Ngoài
ra, việc bổ sung hai chỉ tiêu thường bị bỏ qua -- độ hiệu chuẩn xác suất và thời
gian chấm điểm -- cho phép đánh giá khả năng triển khai chứ không chỉ chất lượng
thống kê.

\subsection{Nhận diện hạn chế và định hướng phát triển}
\label{subsec:han-che-dinh-huong}

\subsubsection{Hạn chế của kết quả thực nghiệm}

\paragraph{Hạn chế về dữ liệu.} Đây là hạn chế cơ bản nhất. Dữ liệu mô phỏng được
sinh theo một mô hình xác suất do chính đề án thiết kế, nên cấu trúc phụ thuộc
giữa các biến phản ánh giả định của người thiết kế chứ không phải quy luật thực
tế. Ba hệ quả cụ thể: (i) giá trị tuyệt đối của các chỉ tiêu chất lượng không có
ý nghĩa dự báo về hiệu quả trên dữ liệu thực; (ii) quan hệ giữa các đặc trưng
thực tế có thể phức tạp hơn nhiều so với mô hình sinh; (iii) dữ liệu thực còn
chứa các dạng nhiễu mà mô hình sinh không tái tạo được, như sai sót nhập liệu có
hệ thống hay hành vi khai báo thiếu trung thực.

\paragraph{Hạn chế về phạm vi đo lường.} Các số đo thời gian được thực hiện trên
một máy đơn, một tiến trình, không bao gồm thời gian truyền qua mạng, truy vấn cơ
sở dữ liệu nguồn hay ảnh hưởng của tải đồng thời. Chúng có giá trị so sánh tương
đối giữa các mô hình nhưng không phải dự báo về thời gian phản hồi trong môi
trường sản xuất.

\paragraph{Hạn chế về phạm vi nghiệp vụ.} Đề án chỉ mô hình hóa thành phần $\PD$
trong phân rã tổn thất kỳ vọng; $\LGD$ và $\EAD$ không thuộc phạm vi. Ngoài ra,
đề án chưa xử lý bài toán \eng{reject inference}: mô hình được huấn luyện trên
các hồ sơ đã được cấp tín dụng, trong khi trong thực tế còn một nhóm hồ sơ bị từ
chối không có kết quả trả nợ. Bỏ qua nhóm này có thể dẫn tới sai lệch chọn mẫu.

\paragraph{Hạn chế về thời gian quan sát.} Đánh giá độ ổn định chỉ dựa trên bốn
tháng dữ liệu kiểm thử -- quá ngắn so với chu kỳ suy giảm điển hình của mô hình
tín dụng. Kết luận về độ ổn định vì vậy mang tính tham khảo.

\paragraph{Hạn chế về kiểm định công bằng.} Đề án chưa thực hiện kiểm định công
bằng thuật toán (\eng{fairness testing}) -- đánh giá xem mô hình có tạo ra chênh
lệch có hệ thống giữa các nhóm khách hàng hay không. Đây là yêu cầu ngày càng
được chú trọng đối với mô hình sử dụng trong cấp tín dụng.

\subsubsection{Định hướng phát triển}

Từ các hạn chế trên, đề án xác định năm hướng phát triển tiếp theo:

\begin{enumerate}[label=(\arabic*)]
  \item \textbf{Kiểm chứng trên dữ liệu thực.} Áp dụng toàn bộ quy trình đã xây
  dựng trên dữ liệu ẩn danh của một tổ chức tín dụng để đối chiếu kết quả. Đây là
  hướng có giá trị cao nhất vì nó trực tiếp giải quyết hạn chế cơ bản nhất.

  \item \textbf{Bổ sung xử lý reject inference.} Áp dụng các phương pháp suy luận
  nhãn cho nhóm hồ sơ bị từ chối để khắc phục sai lệch chọn mẫu.

  \item \textbf{Mở rộng sang mô hình $\LGD$ và $\EAD$.} Hoàn thiện khung tính tổn
  thất kỳ vọng, cho phép định giá theo rủi ro một cách đầy đủ.

  \item \textbf{Bổ sung kiểm định công bằng và giám sát liên tục.} Xây dựng bộ chỉ
  tiêu đánh giá công bằng và cơ chế cảnh báo tự động khi phát hiện suy giảm chất
  lượng hoặc chênh lệch giữa các nhóm khách hàng.

  \item \textbf{Nghiên cứu cơ chế cập nhật mô hình.} Thiết kế quy trình huấn luyện
  lại định kỳ có kiểm soát, bao gồm tiêu chí kích hoạt, quy trình kiểm định mô
  hình mới trước khi thay thế và cơ chế quay lui khi phát hiện vấn đề.
\end{enumerate}

\section*{Kết luận chương 3}
\addcontentsline{toc}{section}{Kết luận chương 3}

Chương 3 đã trình bày toàn bộ kết quả thực nghiệm của đề án, với năm nhóm kết quả
chính.

\textit{Về dữ liệu và đặc trưng}: tập dữ liệu gồm \NTotal{} hồ sơ vay với
\NFeaturesInitial{} đặc trưng ban đầu, rút gọn còn \NFeaturesSelected{} biến sau
ba bộ lọc. Cơ chế kiểm tra chất lượng dữ liệu được kiểm chứng bằng phép tiêm lỗi
với tỷ lệ phát hiện \ValidityDetected{}; phân tích vintage xác nhận cửa sổ quan
sát 12 tháng ghi nhận được \VintageCoverage{} tổng số trường hợp vỡ nợ.

\textit{Về chất lượng mô hình}: \NModels{} mô hình được huấn luyện và so sánh
trên cùng điều kiện. Mô hình \BestModel{} đạt $\AUC = \BestAUC{}$,
$\text{Gini} = \BestGini{}$ và $\KS = \BestKS{}$, còn cách trần Bayes \BayesAUC{}
\GapBayes{} điểm. Sai số hiệu chuẩn kỳ vọng \BestECE{} cho phép sử dụng trực
tiếp giá trị $\PD$ trong phân vùng quyết định.

\textit{Về so sánh thuật toán}: chênh lệch giữa mô hình tốt nhất và hồi quy
logistic chỉ \GapBestLR{} điểm $\AUC$, phù hợp với các công trình đối sánh quốc
tế và đặt ra câu hỏi về việc đánh đổi giữa mức cải thiện nhỏ với chi phí mất khả
năng giải thích.

\textit{Về giải thích mô hình}: nhóm đặc trưng hành vi giao dịch đóng góp
\ShapBehaviourPct{}\% vào quyết định của mô hình, vượt nhóm CIC truyền thống
(\ShapCICPct{}\%).

\textit{Về mô phỏng quy trình kiểm duyệt}: mô phỏng trên tập kiểm thử cho tỷ lệ
xử lý tự động \AutomationRate{} với tỷ lệ nợ xấu thực tế tăng đơn điệu qua ba nhóm khuyến nghị
(\BadApprove{} -- \BadReview{} -- \BadReject{}). So với quy trình thẩm định
truyền thống, mô hình giảm tỷ lệ nợ xấu trong nhóm được duyệt từ \TradBadRate{}
xuống \PropBadRate{} (giảm tương đối \BadRateReductionPct{}), đổi lại giảm
\ApproveRateDrop{} điểm phần trăm tỷ lệ duyệt và tăng \ReviewRateRise{} điểm phần
trăm khối lượng thẩm định thủ công. Thời gian xử lý một hồ sơ ở phân vị 95\% là
\SysPNinetyFive{} ms, đạt mục tiêu đề ra.
